% Options for packages loaded elsewhere
\PassOptionsToPackage{unicode}{hyperref}
\PassOptionsToPackage{hyphens}{url}
\PassOptionsToPackage{dvipsnames,svgnames,x11names}{xcolor}
%
\documentclass[
11pt,
a4paper,
]{report}
\usepackage{amsmath,amssymb}
\usepackage{iftex}
\ifPDFTeX
\usepackage[T1]{fontenc}
\usepackage[utf8]{inputenc}
\usepackage{textcomp} % provide euro and other symbols
\else % if luatex or xetex
\usepackage{unicode-math} % this also loads fontspec
\defaultfontfeatures{Scale=MatchLowercase}
\defaultfontfeatures[\rmfamily]{Ligatures=TeX,Scale=1}
\fi
\usepackage[]{mathpazo}
\ifPDFTeX\else
% xetex/luatex font selection
\fi
% Use upquote if available, for straight quotes in verbatim environments
\IfFileExists{upquote.sty}{\usepackage{upquote}}{}
\IfFileExists{microtype.sty}{% use microtype if available
	\usepackage[]{microtype}
	\UseMicrotypeSet[protrusion]{basicmath} % disable protrusion for tt fonts
}{}
\makeatletter
\@ifundefined{KOMAClassName}{% if non-KOMA class
	\IfFileExists{parskip.sty}{%
		\usepackage{parskip}
	}{% else
		\setlength{\parindent}{0pt}
		\setlength{\parskip}{6pt plus 2pt minus 1pt}}
}{% if KOMA class
	\KOMAoptions{parskip=half}}
\makeatother
\usepackage{xcolor}
\usepackage[margin=2.6cm]{geometry}
\usepackage{listings}
\newcommand{\passthrough}[1]{#1}
\lstset{defaultdialect=[5.3]Lua}
\lstset{defaultdialect=[x86masm]Assembler}
\usepackage{longtable,booktabs,array}
\usepackage{calc} % for calculating minipage widths
% Correct order of tables after \paragraph or \subparagraph
\usepackage{etoolbox}
\makeatletter
\patchcmd\longtable{\par}{\if@noskipsec\mbox{}\fi\par}{}{}
\makeatother
% Allow footnotes in longtable head/foot
\IfFileExists{footnotehyper.sty}{\usepackage{footnotehyper}}{\usepackage{footnote}}
\makesavenoteenv{longtable}
\setlength{\emergencystretch}{3em} % prevent overfull lines
\providecommand{\tightlist}{%
	\setlength{\itemsep}{0pt}\setlength{\parskip}{0pt}}
\setcounter{secnumdepth}{-\maxdimen} % remove section numbering
\ifLuaTeX
\usepackage[bidi=basic]{babel}
\else
\usepackage[bidi=default]{babel}
\fi
\babelprovide[main,import]{spanish}
% get rid of language-specific shorthands (see #6817):
\let\LanguageShortHands\languageshorthands
\def\languageshorthands#1{}
% ==================== Preámbulo de formato ====================
\usepackage{geometry}
\geometry{a4paper, top=2.7cm, bottom=2.7cm, left=2.6cm, right=2.6cm}

\usepackage{xcolor}
\definecolor{primaryblue}{HTML}{1B3A5C}
\definecolor{accentblue}{HTML}{2E6DA4}
\definecolor{codebg}{HTML}{F6F8FA}
\definecolor{codeframe}{HTML}{D0D7DE}
\definecolor{stringgreen}{HTML}{116329}
\definecolor{commentgray}{HTML}{6E7781}
\definecolor{keywordpurple}{HTML}{8250DF}

% ---------- Títulos de sección ----------
\usepackage{titlesec}
\titleformat{\chapter}[display]
{\normalfont\huge\bfseries\color{black}}
{\chaptertitlename\ \thechapter}{18pt}{\Huge}
\titleformat{\section}
{\normalfont\Large\bfseries\color{black}}{\thesection}{1em}{}
\titleformat{\subsection}
{\normalfont\large\bfseries\color{black}}{\thesubsection}{1em}{}
\titleformat{\subsubsection}
{\normalfont\normalsize\bfseries\color{black}}{\thesubsubsection}{1em}{}
\titlespacing*{\chapter}{0pt}{0pt}{24pt}

% ---------- Encabezados y pies de página ----------
\usepackage{fancyhdr}
\setlength{\headheight}{22pt}
\pagestyle{fancy}
\fancyhf{}
\fancyhead[L]{\small\textcolor{gray}{SIGCB-QR --- Tercera Entrega}}
\fancyhead[R]{\small\textcolor{gray}{\leftmark}}
\fancyfoot[C]{\small\thepage}
\renewcommand{\headrulewidth}{0.4pt}
\renewcommand{\headrule}{\hbox to\headwidth{\color{codeframe}\leaders\hrule height \headrulewidth\hfill}}
\fancypagestyle{plain}{%
	\fancyhf{}
	\fancyfoot[C]{\small\thepage}
	\renewcommand{\headrulewidth}{0pt}
}

% ---------- Tablas ----------
\usepackage{booktabs}
\usepackage{longtable}
\usepackage{array}
\usepackage{pdflscape}
\usepackage{seqsplit}
\renewcommand{\arraystretch}{1.25}

% ---------- Código fuente (listings) ----------
\usepackage{listings}
\lstdefinelanguage{gherkin}{
	morekeywords={Feature,Scenario,Scenario Outline,Given,When,Then,And,But,Background,Examples},
	sensitive=true,
	morecomment=[l]{\#},
	morestring=[b]"
}
\lstset{
	extendedchars=true,
	inputencoding=utf8,
	literate=
	{á}{{\'a}}1 {é}{{\'e}}1 {í}{{\'i}}1 {ó}{{\'o}}1 {ú}{{\'u}}1
	{Á}{{\'A}}1 {É}{{\'E}}1 {Í}{{\'I}}1 {Ó}{{\'O}}1 {Ú}{{\'U}}1
	{ñ}{{\~n}}1 {Ñ}{{\~N}}1 {ü}{{\"u}}1 {Ü}{{\"U}}1
	{¿}{{?`}}1 {¡}{{!`}}1,
	backgroundcolor=\color{codebg},
	basicstyle=\ttfamily\small,
	keywordstyle=\color{keywordpurple}\bfseries,
	commentstyle=\color{commentgray}\itshape,
	stringstyle=\color{stringgreen},
	numbers=left,
	numberstyle=\tiny\color{gray},
	numbersep=8pt,
	frame=single,
	rulecolor=\color{codeframe},
	framesep=6pt,
	breaklines=true,
	breakatwhitespace=false,
	showstringspaces=false,
	tabsize=2,
	captionpos=b,
	columns=flexible,
	xleftmargin=1pt,
	xrightmargin=1pt
}

% ---------- Hipervínculos ----------
\usepackage{hyperref}
\hypersetup{
	colorlinks=true,
	linkcolor=accentblue,
	citecolor=accentblue,
	urlcolor=accentblue,
	pdftitle={SIGCB-QR - Sistema Integral de Gestion Bibliotecaria},
	pdfauthor={Equipo SIGCB-QR}
}

% ---------- Otros ----------
\usepackage{enumitem}
\setlist{itemsep=2pt, topsep=4pt}
\usepackage{parskip}
\setlength{\parskip}{6pt}
\ifLuaTeX
\usepackage{selnolig}  % disable illegal ligatures
\fi
\IfFileExists{bookmark.sty}{\usepackage{bookmark}}{\usepackage{hyperref}}
\IfFileExists{xurl.sty}{\usepackage{xurl}}{} % add URL line breaks if available
\urlstyle{same}
\hypersetup{
	pdflang={es-ES},
	colorlinks=true,
	linkcolor={Maroon},
	filecolor={Maroon},
	citecolor={Blue},
	urlcolor={Blue},
	pdfcreator={LaTeX via pandoc}}

\author{}
\date{}

\begin{document}
	
	\begin{titlepage}
		\centering
		\vspace*{2.5cm}
		
		{\Large\bfseries\color{primaryblue} TERCERA ENTREGA PFC}\\[0.6cm]
		
		{\huge\bfseries\color{primaryblue} SIGCB-QR}\\[0.3cm]
		{\Large\color{accentblue} Sistema Integral de Gestión Bibliotecaria}\\[1.5cm]
		
		\rule{0.55\textwidth}{0.6pt}\\[0.5cm]
		
		{\large Documento consolidado: SRS, Casos de Uso, Historias de Usuario,\\
			Matriz de Trazabilidad, Código Fuente y Apéndices}\\[2cm]
		
		\begin{tabular}{rl}
			\textbf{Versión:} & v0.9.0-rc \\[4pt]
			\textbf{Fecha:} & Julio 2026 \\[4pt]
			\textbf{Conforme a:} & ISO/IEC/IEEE 29148:2018 \\[4pt]
			\textbf{Equipo:} & ARZ \\
			\textbf{ } & \shortstack[l]{ARIAS MOREIRA MAYBELIN GREGORIA \\ ROMERO MENDEZ BRYAM STEVEN \\ ZAMBRANO MOREIRA ALISON ARIANA} \\
		\end{tabular}
		
		\vfill
		
		{\small\color{gray} Documento técnico generado a partir del material consolidado del proyecto}
		
	\end{titlepage}
	
	{
		\hypersetup{linkcolor=}
		\setcounter{tocdepth}{3}
		\tableofcontents
	}
	\hypertarget{software-requirements-specification-srs}{%
		\chapter{Software Requirements Specification (SRS)}\label{software-requirements-specification-srs}}
	
	\textbf{Version:} v0.9.0-rc
	\textbf{Fecha:} 2026-07-30
	\textbf{Conforme a:} ISO/IEC/IEEE 29148:2018
	
	\hypertarget{introduccion}{%
		\section{1. Introduccion}\label{introduccion}}
	
	\hypertarget{proposito}{%
		\subsection{1.1 Proposito}\label{proposito}}
	
	Este documento especifica los requisitos funcionales y no funcionales del Sistema Integral de Gestion, Control y Seguimiento de Lectura Interna en la Biblioteca Universitaria mediante Codigos QR (SIGCB-QR). El sistema automatiza los procesos bibliotecarios: autenticacion, gestion de catalogo, prestamos, reservas, multas y reportes.
	
	\hypertarget{alcance}{%
		\subsection{1.2 Alcance}\label{alcance}}
	
	SIGCB-QR es una aplicacion web de arquitectura cliente-servidor con backend Spring Boot 3.4, frontend Laravel 13, base de datos PostgreSQL 16 y cache Redis 7. El sistema expone una API REST documentada con OpenAPI 3.0 y utiliza autenticacion stateless JWT en cookie HttpOnly.
	
	\hypertarget{definiciones}{%
		\subsection{1.3 Definiciones}\label{definiciones}}
	
	\begin{longtable}[]{@{}ll@{}}
		\toprule\noalign{}
		Termino & Definicion \\
		\midrule\noalign{}
		\endhead
		\bottomrule\noalign{}
		\endlastfoot
		JWT & JSON Web Token (RFC 7519) \\
		JPA & Jakarta Persistence API \\
		ORM & Object-Relational Mapping \\
		CRUD & Create, Read, Update, Delete \\
		TTL & Time To Live \\
		Hit ratio & Proporcion de aciertos de cache \\
	\end{longtable}
	
	\hypertarget{referencias}{%
		\subsection{1.4 Referencias}\label{referencias}}
	
	\begin{itemize}
		\tightlist
		\item
		ISO/IEC/IEEE 29148:2018 -- Ingenieria de requisitos
		\item
		RFC 7519 -- JSON Web Token
		\item
		RFC 7807 -- Problem Details for HTTP APIs
		\item
		INCOSE Guide to Writing Requirements v4
		\item
		OWASP Cheat Sheet -- SQL Injection Prevention
		\item
		Cockburn -- Writing Effective Use Cases
	\end{itemize}
	
	\hypertarget{resumen}{%
		\subsection{1.5 Resumen}\label{resumen}}
	
	El SRS se organiza en: descripcion global (seccion 2), requisitos especificos funcionales (3), no funcionales (4), de interfaz externa (5), de seguridad (6) y de calidad (7).
	
	\hypertarget{descripcion-global}{%
		\section{2. Descripcion global}\label{descripcion-global}}
	
	\hypertarget{perspectiva-del-producto}{%
		\subsection{2.1 Perspectiva del producto}\label{perspectiva-del-producto}}
	
	SIGCB-QR es un sistema independiente que reemplaza procesos manuales bibliotecarios. Se compone de:
	- API REST (sigcb-qr-api): Java 21, Spring Boot 3.4.4
	- Frontend (SIGCB-QR): Laravel 13, PHP 8.4
	- Base de datos: PostgreSQL 16
	- Cache: Redis 7
	- Contenedores: Docker Compose
	
	\hypertarget{funciones-del-sistema}{%
		\subsection{2.2 Funciones del sistema}\label{funciones-del-sistema}}
	
	\begin{itemize}
		\tightlist
		\item
		Autenticacion stateless con JWT (login, registro, logout)
		\item
		CRUD de usuarios (admin)
		\item
		CRUD de material bibliografico (libros, autores, editoriales, categorias)
		\item
		Gestion de ejemplares (inventario)
		\item
		Prestamos, devoluciones y renovaciones
		\item
		Reservas
		\item
		Multas y sanciones
		\item
		Reportes y estadisticas de dashboard
		\item
		Catalogacion (facultades, carreras)
	\end{itemize}
	
	\hypertarget{caracteristicas-de-usuarios}{%
		\subsection{2.3 Caracteristicas de usuarios}\label{caracteristicas-de-usuarios}}
	
	\begin{longtable}[]{@{}ll@{}}
		\toprule\noalign{}
		Rol & Descripcion \\
		\midrule\noalign{}
		\endhead
		\bottomrule\noalign{}
		\endlastfoot
		ADMIN & Administracion completa del sistema \\
		BIBLIOTECARIO & Gestion de prestamos, catalogo y reportes \\
		ESTUDIANTE & Consulta de catalogo, prestamos y reservas \\
	\end{longtable}
	
	\hypertarget{restricciones}{%
		\subsection{2.4 Restricciones}\label{restricciones}}
	
	\begin{itemize}
		\tightlist
		\item
		Autenticacion stateless (sin sesiones HTTP)
		\item
		JWT transmitido en cookie HttpOnly + Secure + SameSite=Strict
		\item
		Base de datos gestionada exclusivamente por Flyway (ddl-auto=validate)
		\item
		Operaciones no elementales mediante stored procedures
		\item
		Contenedores Docker anclados por digest SHA256
	\end{itemize}
	
	\hypertarget{supuestos-y-dependencias}{%
		\subsection{2.5 Supuestos y dependencias}\label{supuestos-y-dependencias}}
	
	\begin{itemize}
		\tightlist
		\item
		Disponibilidad de Docker y Docker Compose
		\item
		Conexion de red entre contenedores
		\item
		Semillas deterministas fijadas
	\end{itemize}
	
	\hypertarget{requisitos-especificos-funcionales}{%
		\section{3. Requisitos especificos -- Funcionales}\label{requisitos-especificos-funcionales}}
	
	\hypertarget{req-f-001-autenticacion-de-usuarios}{%
		\subsection{REQ-F-001 -- Autenticacion de usuarios}\label{req-f-001-autenticacion-de-usuarios}}
	
	\begin{longtable}[]{@{}
			>{\raggedright\arraybackslash}p{(\columnwidth - 2\tabcolsep) * \real{0.5882}}
			>{\raggedright\arraybackslash}p{(\columnwidth - 2\tabcolsep) * \real{0.4118}}@{}}
		\toprule\noalign{}
		\begin{minipage}[b]{\linewidth}\raggedright
			Atributo
		\end{minipage} & \begin{minipage}[b]{\linewidth}\raggedright
			Valor
		\end{minipage} \\
		\midrule\noalign{}
		\endhead
		\bottomrule\noalign{}
		\endlastfoot
		Tipo & Funcional \\
		Prioridad (MoSCoW) & Must \\
		Enunciado & Al recibir credenciales validas, el sistema debera emitir un JWT firmado en cookie HttpOnly con los claims estandar (iss, sub, aud, exp, nbf, iat, jti) \\
		Rationale & La autenticacion stateless elimina la necesidad de almacenar sesiones en servidor \\
		Origen & Stakeholder: Administracion bibliotecaria \\
		Criterio de aceptacion & Login exitoso en menos de 2s con JWT valido \\
		Metodo de verificacion & Test (AuthControllerTest) \\
		Trazabilidad & HU-01, CU-01 \\
	\end{longtable}
	
	\hypertarget{req-f-002-registro-de-usuarios}{%
		\subsection{REQ-F-002 -- Registro de usuarios}\label{req-f-002-registro-de-usuarios}}
	
	\begin{longtable}[]{@{}
			>{\raggedright\arraybackslash}p{(\columnwidth - 2\tabcolsep) * \real{0.5882}}
			>{\raggedright\arraybackslash}p{(\columnwidth - 2\tabcolsep) * \real{0.4118}}@{}}
		\toprule\noalign{}
		\begin{minipage}[b]{\linewidth}\raggedright
			Atributo
		\end{minipage} & \begin{minipage}[b]{\linewidth}\raggedright
			Valor
		\end{minipage} \\
		\midrule\noalign{}
		\endhead
		\bottomrule\noalign{}
		\endlastfoot
		Tipo & Funcional \\
		Prioridad & Must \\
		Enunciado & Al recibir datos de registro validos, el sistema debera crear un usuario con rol ESTUDIANTE por defecto y emitir un JWT \\
		Rationale & Permitir auto-registro de estudiantes \\
		Criterio de aceptacion & Registro exitoso con validacion de email unico \\
		Metodo de verificacion & Test \\
		Trazabilidad & HU-02, CU-01 \\
	\end{longtable}
	
	\hypertarget{req-f-003-cierre-de-sesion}{%
		\subsection{REQ-F-003 -- Cierre de sesion}\label{req-f-003-cierre-de-sesion}}
	
	\begin{longtable}[]{@{}
			>{\raggedright\arraybackslash}p{(\columnwidth - 2\tabcolsep) * \real{0.5882}}
			>{\raggedright\arraybackslash}p{(\columnwidth - 2\tabcolsep) * \real{0.4118}}@{}}
		\toprule\noalign{}
		\begin{minipage}[b]{\linewidth}\raggedright
			Atributo
		\end{minipage} & \begin{minipage}[b]{\linewidth}\raggedright
			Valor
		\end{minipage} \\
		\midrule\noalign{}
		\endhead
		\bottomrule\noalign{}
		\endlastfoot
		Tipo & Funcional \\
		Prioridad & Must \\
		Enunciado & Al recibir una solicitud de logout, el sistema debera agregar el JTI del token a la blacklist Redis e invalidar la cookie \\
		Rationale & Garantizar que el token no pueda reutilizarse \\
		Criterio de aceptacion & Token blacklistado no es valido para peticiones posteriores \\
		Metodo de verificacion & Test \\
		Trazabilidad & HU-03, CU-01 \\
	\end{longtable}
	
	\hypertarget{req-f-004-crud-de-usuarios-admin}{%
		\subsection{REQ-F-004 -- CRUD de usuarios (Admin)}\label{req-f-004-crud-de-usuarios-admin}}
	
	\begin{longtable}[]{@{}
			>{\raggedright\arraybackslash}p{(\columnwidth - 2\tabcolsep) * \real{0.5882}}
			>{\raggedright\arraybackslash}p{(\columnwidth - 2\tabcolsep) * \real{0.4118}}@{}}
		\toprule\noalign{}
		\begin{minipage}[b]{\linewidth}\raggedright
			Atributo
		\end{minipage} & \begin{minipage}[b]{\linewidth}\raggedright
			Valor
		\end{minipage} \\
		\midrule\noalign{}
		\endhead
		\bottomrule\noalign{}
		\endlastfoot
		Tipo & Funcional \\
		Prioridad & Must \\
		Enunciado & El sistema debera permitir al administrador crear, leer, actualizar y desactivar usuarios \\
		Rationale & Gestion centralizada de cuentas \\
		Criterio de aceptacion & CRUD completo con validaciones \\
		Metodo de verificacion & Test \\
		Trazabilidad & HU-04 \\
	\end{longtable}
	
	\hypertarget{req-f-005-crud-de-material-bibliografico}{%
		\subsection{REQ-F-005 -- CRUD de material bibliografico}\label{req-f-005-crud-de-material-bibliografico}}
	
	\begin{longtable}[]{@{}
			>{\raggedright\arraybackslash}p{(\columnwidth - 2\tabcolsep) * \real{0.5882}}
			>{\raggedright\arraybackslash}p{(\columnwidth - 2\tabcolsep) * \real{0.4118}}@{}}
		\toprule\noalign{}
		\begin{minipage}[b]{\linewidth}\raggedright
			Atributo
		\end{minipage} & \begin{minipage}[b]{\linewidth}\raggedright
			Valor
		\end{minipage} \\
		\midrule\noalign{}
		\endhead
		\bottomrule\noalign{}
		\endlastfoot
		Tipo & Funcional \\
		Prioridad & Must \\
		Enunciado & El sistema debera permitir crear, leer, actualizar y desactivar libros con sus metadatos (titulo, ISBN, autor, editorial, categoria) \\
		Rationale & Catalogacion del acervo bibliotecario \\
		Criterio de aceptacion & CRUD completo con paginacion y busqueda \\
		Metodo de verificacion & Test \\
		Trazabilidad & HU-05, CU-02 \\
	\end{longtable}
	
	\hypertarget{req-f-006-prestamo-de-ejemplares}{%
		\subsection{REQ-F-006 -- Prestamo de ejemplares}\label{req-f-006-prestamo-de-ejemplares}}
	
	\begin{longtable}[]{@{}
			>{\raggedright\arraybackslash}p{(\columnwidth - 2\tabcolsep) * \real{0.5882}}
			>{\raggedright\arraybackslash}p{(\columnwidth - 2\tabcolsep) * \real{0.4118}}@{}}
		\toprule\noalign{}
		\begin{minipage}[b]{\linewidth}\raggedright
			Atributo
		\end{minipage} & \begin{minipage}[b]{\linewidth}\raggedright
			Valor
		\end{minipage} \\
		\midrule\noalign{}
		\endhead
		\bottomrule\noalign{}
		\endlastfoot
		Tipo & Funcional \\
		Prioridad & Must \\
		Enunciado & Al recibir una solicitud de prestamo valida, el sistema debera crear el prestamo, actualizar el estado del ejemplar y disminuir los ejemplares disponibles \\
		Rationale & Control de circulacion de material \\
		Criterio de aceptacion & Prestamo creado en menos de 2s, maximo 5 activos por usuario \\
		Metodo de verificacion & Test + stored procedure \\
		Trazabilidad & HU-06, CU-03 \\
	\end{longtable}
	
	\hypertarget{req-f-007-devolucion-de-prestamo}{%
		\subsection{REQ-F-007 -- Devolucion de prestamo}\label{req-f-007-devolucion-de-prestamo}}
	
	\begin{longtable}[]{@{}
			>{\raggedright\arraybackslash}p{(\columnwidth - 2\tabcolsep) * \real{0.5882}}
			>{\raggedright\arraybackslash}p{(\columnwidth - 2\tabcolsep) * \real{0.4118}}@{}}
		\toprule\noalign{}
		\begin{minipage}[b]{\linewidth}\raggedright
			Atributo
		\end{minipage} & \begin{minipage}[b]{\linewidth}\raggedright
			Valor
		\end{minipage} \\
		\midrule\noalign{}
		\endhead
		\bottomrule\noalign{}
		\endlastfoot
		Tipo & Funcional \\
		Prioridad & Must \\
		Enunciado & Al recibir una devolucion, el sistema debera actualizar el estado del prestamo, liberar el ejemplar y generar multa si esta vencido \\
		Rationale & Control de inventario y penalizaciones \\
		Criterio de aceptacion & Devolucion procesada en menos de 2s \\
		Metodo de verificacion & Test + stored procedure \\
		Trazabilidad & HU-07, CU-04 \\
	\end{longtable}
	
	\hypertarget{req-f-008-reporte-de-prestamos-diarios}{%
		\subsection{REQ-F-008 -- Reporte de prestamos diarios}\label{req-f-008-reporte-de-prestamos-diarios}}
	
	\begin{longtable}[]{@{}
			>{\raggedright\arraybackslash}p{(\columnwidth - 2\tabcolsep) * \real{0.5882}}
			>{\raggedright\arraybackslash}p{(\columnwidth - 2\tabcolsep) * \real{0.4118}}@{}}
		\toprule\noalign{}
		\begin{minipage}[b]{\linewidth}\raggedright
			Atributo
		\end{minipage} & \begin{minipage}[b]{\linewidth}\raggedright
			Valor
		\end{minipage} \\
		\midrule\noalign{}
		\endhead
		\bottomrule\noalign{}
		\endlastfoot
		Tipo & Funcional \\
		Prioridad & Should \\
		Enunciado & El sistema debera generar un reporte de prestamos del dia actual con datos de usuario y libro \\
		Rationale & Seguimiento de actividad diaria \\
		Criterio de aceptacion & Reporte generado en menos de 3s \\
		Metodo de verificacion & Test + stored procedure \\
		Trazabilidad & HU-08 \\
	\end{longtable}
	
	\hypertarget{requisitos-no-funcionales}{%
		\section{4. Requisitos no funcionales}\label{requisitos-no-funcionales}}
	
	\hypertarget{req-nf-001-rendimiento-listado-de-libros}{%
		\subsection{REQ-NF-001 -- Rendimiento -- Listado de libros}\label{req-nf-001-rendimiento-listado-de-libros}}
	
	\begin{longtable}[]{@{}
			>{\raggedright\arraybackslash}p{(\columnwidth - 2\tabcolsep) * \real{0.5882}}
			>{\raggedright\arraybackslash}p{(\columnwidth - 2\tabcolsep) * \real{0.4118}}@{}}
		\toprule\noalign{}
		\begin{minipage}[b]{\linewidth}\raggedright
			Atributo
		\end{minipage} & \begin{minipage}[b]{\linewidth}\raggedright
			Valor
		\end{minipage} \\
		\midrule\noalign{}
		\endhead
		\bottomrule\noalign{}
		\endlastfoot
		Tipo & No funcional \\
		Prioridad & Must \\
		Enunciado & El endpoint GET /api/libros debera responder con p95 \textless= 200 ms con cache Redis caliente \\
		Rationale & Experiencia de usuario fluida en navegacion de catalogo \\
		Criterio de aceptacion & p95 \textless{} 200ms en 3 corridas con cache caliente \\
		Metodo de verificacion & Test de rendimiento (k6) \\
		Trazabilidad & Evidencia: docs/mediciones/perf/ \\
	\end{longtable}
	
	\hypertarget{req-nf-002-seguridad-autenticacion}{%
		\subsection{REQ-NF-002 -- Seguridad -- Autenticacion}\label{req-nf-002-seguridad-autenticacion}}
	
	\begin{longtable}[]{@{}
			>{\raggedright\arraybackslash}p{(\columnwidth - 2\tabcolsep) * \real{0.5882}}
			>{\raggedright\arraybackslash}p{(\columnwidth - 2\tabcolsep) * \real{0.4118}}@{}}
		\toprule\noalign{}
		\begin{minipage}[b]{\linewidth}\raggedright
			Atributo
		\end{minipage} & \begin{minipage}[b]{\linewidth}\raggedright
			Valor
		\end{minipage} \\
		\midrule\noalign{}
		\endhead
		\bottomrule\noalign{}
		\endlastfoot
		Tipo & No funcional \\
		Prioridad & Must \\
		Enunciado & El sistema debera implementar autenticacion JWT sin sesiones, con cookie HttpOnly + Secure + SameSite=Strict \\
		Rationale & Mitigacion de ataques XSS y CSRF \\
		Criterio de aceptacion & Ningun endpoint protegido accesible sin JWT valido \\
		Metodo de verificacion & Test de seguridad \\
		Trazabilidad & OWASP Top 10 \\
	\end{longtable}
	
	\hypertarget{req-nf-003-seguridad-proteccion-contra-inyeccion-sql}{%
		\subsection{REQ-NF-003 -- Seguridad -- Proteccion contra inyeccion SQL}\label{req-nf-003-seguridad-proteccion-contra-inyeccion-sql}}
	
	\begin{longtable}[]{@{}
			>{\raggedright\arraybackslash}p{(\columnwidth - 2\tabcolsep) * \real{0.5882}}
			>{\raggedright\arraybackslash}p{(\columnwidth - 2\tabcolsep) * \real{0.4118}}@{}}
		\toprule\noalign{}
		\begin{minipage}[b]{\linewidth}\raggedright
			Atributo
		\end{minipage} & \begin{minipage}[b]{\linewidth}\raggedright
			Valor
		\end{minipage} \\
		\midrule\noalign{}
		\endhead
		\bottomrule\noalign{}
		\endlastfoot
		Tipo & No funcional \\
		Prioridad & Must \\
		Enunciado & Ninguna consulta SQL debera concatenar entrada de usuario; toda parametrizacion sera mediante parametros nombrados \\
		Rationale & Prevencion de inyeccion SQL \\
		Criterio de aceptacion & Escaneo SAST sin hallazgos de inyeccion SQL \\
		Metodo de verificacion & Analisis estatico (inspection) \\
		Trazabilidad & OWASP Cheat Sheet \\
	\end{longtable}
	
	\hypertarget{req-nf-004-disponibilidad-contenedores}{%
		\subsection{REQ-NF-004 -- Disponibilidad -- Contenedores}\label{req-nf-004-disponibilidad-contenedores}}
	
	\begin{longtable}[]{@{}
			>{\raggedright\arraybackslash}p{(\columnwidth - 2\tabcolsep) * \real{0.5882}}
			>{\raggedright\arraybackslash}p{(\columnwidth - 2\tabcolsep) * \real{0.4118}}@{}}
		\toprule\noalign{}
		\begin{minipage}[b]{\linewidth}\raggedright
			Atributo
		\end{minipage} & \begin{minipage}[b]{\linewidth}\raggedright
			Valor
		\end{minipage} \\
		\midrule\noalign{}
		\endhead
		\bottomrule\noalign{}
		\endlastfoot
		Tipo & No funcional \\
		Prioridad & Must \\
		Enunciado & El sistema debera poder reconstruirse en menos de 2 minutos desde una clonacion limpia con un solo comando (make up) \\
		Rationale & Reproducibilidad por terceros \\
		Criterio de aceptacion & docker-compose up exitoso en ambiente limpio \\
		Metodo de verificacion & Demostracion (demonstration) \\
		Trazabilidad & Docker Compose \\
	\end{longtable}
	
	\hypertarget{req-nf-005-mantenibilidad-cobertura-de-pruebas}{%
		\subsection{REQ-NF-005 -- Mantenibilidad -- Cobertura de pruebas}\label{req-nf-005-mantenibilidad-cobertura-de-pruebas}}
	
	\begin{longtable}[]{@{}
			>{\raggedright\arraybackslash}p{(\columnwidth - 2\tabcolsep) * \real{0.5882}}
			>{\raggedright\arraybackslash}p{(\columnwidth - 2\tabcolsep) * \real{0.4118}}@{}}
		\toprule\noalign{}
		\begin{minipage}[b]{\linewidth}\raggedright
			Atributo
		\end{minipage} & \begin{minipage}[b]{\linewidth}\raggedright
			Valor
		\end{minipage} \\
		\midrule\noalign{}
		\endhead
		\bottomrule\noalign{}
		\endlastfoot
		Tipo & No funcional \\
		Prioridad & Should \\
		Enunciado & El backend debera tener cobertura de linea \textgreater= 30\% medida por JaCoCo \\
		Rationale & Garantia minima de calidad \\
		Criterio de aceptacion & Reporte JaCoCo \textgreater= 30\% \\
		Metodo de verificacion & Test (JaCoCo) \\
		Trazabilidad & pom.xml \\
	\end{longtable}
	
	\hypertarget{requisitos-de-interfaz-externa}{%
		\section{5. Requisitos de interfaz externa}\label{requisitos-de-interfaz-externa}}
	
	\hypertarget{interfaz-de-usuario}{%
		\subsection{5.1 Interfaz de usuario}\label{interfaz-de-usuario}}
	
	Frontend Laravel 13 con Tailwind CSS, accesible via navegador en puerto 8000.
	
	\hypertarget{interfaz-api-rest}{%
		\subsection{5.2 Interfaz API REST}\label{interfaz-api-rest}}
	
	Documentacion OpenAPI 3.0 en \passthrough{\lstinline!/api-docs!} e interfaz Swagger UI en \passthrough{\lstinline!/swagger-ui.html!}.
	
	\hypertarget{interfaz-de-base-de-datos}{%
		\subsection{5.3 Interfaz de base de datos}\label{interfaz-de-base-de-datos}}
	
	Conexion JDBC a PostgreSQL 16, esquema gestionado por Flyway.
	
	\hypertarget{requisitos-de-seguridad}{%
		\section{6. Requisitos de seguridad}\label{requisitos-de-seguridad}}
	
	\begin{itemize}
		\tightlist
		\item
		JWT con claims estandar RFC 7519
		\item
		Cookie HttpOnly + Secure + SameSite=Strict
		\item
		Blacklist de JTI revocados en Redis
		\item
		CORS restringido a origenes configurados
		\item
		BCrypt para almacenamiento de contrasenas
		\item
		Parametros nombrados en todas las consultas
		\item
		Procedimientos almacenados sin SQL dinamico
	\end{itemize}
	
	\hypertarget{calidad-de-software}{%
		\section{7. Calidad de software}\label{calidad-de-software}}
	
	\hypertarget{caracteristicas-del-conjunto-de-requisitos-incose-c10-c15}{%
		\subsection{7.1 Caracteristicas del conjunto de requisitos (INCOSE C10-C15)}\label{caracteristicas-del-conjunto-de-requisitos-incose-c10-c15}}
	
	\begin{longtable}[]{@{}
			>{\raggedright\arraybackslash}p{(\columnwidth - 2\tabcolsep) * \real{0.4348}}
			>{\raggedright\arraybackslash}p{(\columnwidth - 2\tabcolsep) * \real{0.5652}}@{}}
		\toprule\noalign{}
		\begin{minipage}[b]{\linewidth}\raggedright
			Atributo
		\end{minipage} & \begin{minipage}[b]{\linewidth}\raggedright
			Cumplimiento
		\end{minipage} \\
		\midrule\noalign{}
		\endhead
		\bottomrule\noalign{}
		\endlastfoot
		Complete (C10) & Todos los escenarios principales cubiertos \\
		Consistent (C11) & Sin contradicciones entre requisitos \\
		Feasible (C12) & Implementable con la pila tecnologica actual \\
		Comprehensible (C13) & Redactado en lenguaje claro \\
		Able to be validated (C14) & Cada requisito tiene metodo de verificacion \\
		Correct (C15) & Aprobado por stakeholders \\
	\end{longtable}
	
	\emph{Fin del SRS}
	
	\hypertarget{casos-de-uso}{%
		\chapter{2. Casos de Uso}\label{casos-de-uso}}
	
	\hypertarget{cu-01-autenticaciuxf3n-de-usuarios}{%
		\section{CU-01: Autenticación de usuarios}\label{cu-01-autenticaciuxf3n-de-usuarios}}
	
	\textbf{Nivel:} 1 (Resumen)
	\textbf{Actor principal:} Usuario no autenticado
	\textbf{Stakeholders:} Administración bibliotecaria (desea control de acceso)
	
	\hypertarget{escenario-principal-de-uxe9xito}{%
		\subsection{Escenario principal de éxito}\label{escenario-principal-de-uxe9xito}}
	
	\begin{enumerate}
		\def\labelenumi{\arabic{enumi}.}
		\tightlist
		\item
		El usuario envía sus credenciales (email, contraseña) al sistema
		\item
		El sistema valida las credenciales contra la base de datos
		\item
		El sistema genera un JWT firmado con claims estándar (iss, sub, aud, exp, nbf, iat, jti)
		\item
		El sistema establece una cookie HttpOnly + Secure + SameSite=Strict con el JWT
		\item
		El usuario queda autenticado y puede acceder a recursos protegidos
		\item
		El sistema permite al usuario cerrar sesión, agregando el JTI a la blacklist
	\end{enumerate}
	
	\hypertarget{extensiones}{%
		\subsection{Extensiones}\label{extensiones}}
	
	\hypertarget{a.-credenciales-invuxe1lidas}{%
		\subsubsection{1a. Credenciales inválidas}\label{a.-credenciales-invuxe1lidas}}
	
	1a1. El sistema responde con HTTP 401 y ProblemDetail
	
	\hypertarget{a.-token-expirado}{%
		\subsubsection{2a. Token expirado}\label{a.-token-expirado}}
	
	2a1. El sistema rechaza la petición con HTTP 401
	
	\hypertarget{a.-token-blacklistado}{%
		\subsubsection{3a. Token blacklistado}\label{a.-token-blacklistado}}
	
	3a1. El sistema rechaza la petición con HTTP 401
	
	\hypertarget{requisitos-asociados}{%
		\subsection{Requisitos asociados}\label{requisitos-asociados}}
	
	\begin{itemize}
		\tightlist
		\item
		REQ-F-001 (Login)
		\item
		REQ-F-002 (Registro)
		\item
		REQ-F-003 (Logout)
		\item
		REQ-NF-002 (Seguridad)
	\end{itemize}
	
	\hypertarget{cu-02-gestiuxf3n-de-material-bibliogruxe1fico}{%
		\section{CU-02: Gestión de material bibliográfico}\label{cu-02-gestiuxf3n-de-material-bibliogruxe1fico}}
	
	\textbf{Nivel:} 2 (Etapa de objetivo de usuario)
	\textbf{Actor principal:} Bibliotecario
	\textbf{Precondición:} Autenticado con rol BIBLIOTECARIO o ADMIN
	
	\hypertarget{escenario-principal-de-uxe9xito-1}{%
		\subsection{Escenario principal de éxito}\label{escenario-principal-de-uxe9xito-1}}
	
	\begin{enumerate}
		\def\labelenumi{\arabic{enumi}.}
		\tightlist
		\item
		El bibliotecario solicita el listado de libros (paginado)
		\item
		El sistema retorna los libros activos con paginación, desde caché Redis si está caliente
		\item
		El bibliotecario busca un libro por título
		\item
		El sistema retorna resultados de búsqueda
		\item
		El bibliotecario crea un nuevo libro con título, ISBN, editorial, categoría y autores
		\item
		El sistema valida los datos y crea el libro con ejemplares disponibles = totales
		\item
		El bibliotecario actualiza datos de un libro existente
		\item
		El bibliotecario desactiva un libro (eliminación lógica)
	\end{enumerate}
	
	\hypertarget{extensiones-1}{%
		\subsection{Extensiones}\label{extensiones-1}}
	
	\hypertarget{a.-datos-invuxe1lidos}{%
		\subsubsection{5a. Datos inválidos}\label{a.-datos-invuxe1lidos}}
	
	5a1. El sistema responde con HTTP 400 y errores de validación
	
	\hypertarget{a.-libro-no-encontrado}{%
		\subsubsection{7a. Libro no encontrado}\label{a.-libro-no-encontrado}}
	
	7a1. El sistema responde con HTTP 404 y ProblemDetail
	
	\hypertarget{requisitos-asociados-1}{%
		\subsection{Requisitos asociados}\label{requisitos-asociados-1}}
	
	\begin{itemize}
		\tightlist
		\item
		REQ-F-005 (CRUD libros)
		\item
		REQ-NF-001 (Rendimiento con caché)
	\end{itemize}
	
	\hypertarget{cu-03-prestamo-de-ejemplares}{%
		\section{CU-03: Prestamo de ejemplares}\label{cu-03-prestamo-de-ejemplares}}
	
	\textbf{Nivel:} 2 (Etapa de objetivo de usuario)
	\textbf{Actor principal:} Bibliotecario
	\textbf{Precondicion:} Autenticado con rol BIBLIOTECARIO o ADMIN
	\textbf{Postcondicion:} Prestamo registrado, ejemplar marcado como PRESTADO
	
	\hypertarget{escenario-principal-de-exito}{%
		\subsection{Escenario principal de exito}\label{escenario-principal-de-exito}}
	
	\begin{enumerate}
		\def\labelenumi{\arabic{enumi}.}
		\tightlist
		\item
		El bibliotecario selecciona un usuario del sistema
		\item
		El bibliotecario selecciona un ejemplar disponible del inventario
		\item
		El sistema valida que el usuario no tenga mas de 5 prestamos activos
		\item
		El sistema valida que el ejemplar este DISPONIBLE
		\item
		El sistema crea el registro de prestamo con fecha actual y vencimiento en 7 dias
		\item
		El sistema cambia el estado del ejemplar a PRESTADO
		\item
		El sistema disminuye el contador de ejemplares disponibles del libro
		\item
		El sistema retorna el detalle del prestamo creado
	\end{enumerate}
	
	\hypertarget{extensiones-2}{%
		\subsection{Extensiones}\label{extensiones-2}}
	
	\hypertarget{a.-limite-de-prestamos-alcanzado}{%
		\subsubsection{3a. Limite de prestamos alcanzado}\label{a.-limite-de-prestamos-alcanzado}}
	
	3a1. El sistema rechaza la operacion con HTTP 400
	3a2. Mensaje: ``El usuario tiene demasiados prestamos activos''
	
	\hypertarget{a.-ejemplar-no-disponible}{%
		\subsubsection{4a. Ejemplar no disponible}\label{a.-ejemplar-no-disponible}}
	
	4a1. El sistema rechaza la operacion con HTTP 400
	4a2. Mensaje: ``El ejemplar no esta disponible''
	
	\hypertarget{a.-error-de-base-de-datos}{%
		\subsubsection{5a. Error de base de datos}\label{a.-error-de-base-de-datos}}
	
	5a1. El sistema revierte la transaccion
	5a2. El sistema retorna HTTP 500 con ProblemDetail
	
	\hypertarget{requisitos-asociados-2}{%
		\subsection{Requisitos asociados}\label{requisitos-asociados-2}}
	
	\begin{itemize}
		\tightlist
		\item
		REQ-F-006 (Prestamo de ejemplares)
		\item
		REQ-NF-006 (Control de acceso por roles)
		\item
		REQ-R-003 (Tiempo de respuesta de SP)
	\end{itemize}
	
	\hypertarget{cu-04-devolucion-de-prestamo}{%
		\section{CU-04: Devolucion de prestamo}\label{cu-04-devolucion-de-prestamo}}
	
	\textbf{Nivel:} 2 (Etapa de objetivo de usuario)
	\textbf{Actor principal:} Bibliotecario
	\textbf{Precondicion:} Autenticado con rol BIBLIOTECARIO o ADMIN
	\textbf{Postcondicion:} Prestamo cerrado, inventario actualizado, multa generada si aplica
	
	\hypertarget{escenario-principal-de-exito-1}{%
		\subsection{Escenario principal de exito}\label{escenario-principal-de-exito-1}}
	
	\begin{enumerate}
		\def\labelenumi{\arabic{enumi}.}
		\tightlist
		\item
		El bibliotecario localiza el prestamo activo por su ID
		\item
		El sistema verifica que el prestamo exista y este ACTIVO
		\item
		El sistema registra la fecha de devolucion (fecha actual)
		\item
		El sistema cambia el estado del prestamo a DEVUELTO
		\item
		El sistema cambia el estado del ejemplar a DISPONIBLE
		\item
		El sistema incrementa el contador de ejemplares disponibles del libro
		\item
		Si la fecha actual supera la fecha de vencimiento, el sistema genera una multa
		\item
		El sistema retorna el detalle del prestamo actualizado
	\end{enumerate}
	
	\hypertarget{extensiones-3}{%
		\subsection{Extensiones}\label{extensiones-3}}
	
	\hypertarget{a.-prestamo-no-encontrado}{%
		\subsubsection{2a. Prestamo no encontrado}\label{a.-prestamo-no-encontrado}}
	
	2a1. El sistema retorna HTTP 404 con ProblemDetail
	
	\hypertarget{b.-prestamo-ya-devuelto}{%
		\subsubsection{2b. Prestamo ya devuelto}\label{b.-prestamo-ya-devuelto}}
	
	2b1. El sistema rechaza con HTTP 400
	2b2. Mensaje: ``El prestamo ya fue devuelto''
	
	\hypertarget{a.-calculo-de-multa}{%
		\subsubsection{7a. Calculo de multa}\label{a.-calculo-de-multa}}
	
	7a1. Dias de retraso = fecha\_actual - fecha\_vencimiento
	7a2. Monto = dias\_retraso * \$0.50
	7a3. La multa se registra con estado PENDIENTE
	
	\hypertarget{requisitos-asociados-3}{%
		\subsection{Requisitos asociados}\label{requisitos-asociados-3}}
	
	\begin{itemize}
		\tightlist
		\item
		REQ-F-007 (Devolucion de prestamo)
		\item
		REQ-R-003 (Tiempo de respuesta de SP)
	\end{itemize}
	
	\hypertarget{cu-05-generacion-de-reportes}{%
		\section{CU-05: Generacion de reportes}\label{cu-05-generacion-de-reportes}}
	
	\textbf{Nivel:} 1 (Resumen)
	\textbf{Actor principal:} Bibliotecario/Administrador
	\textbf{Precondicion:} Autenticado con rol BIBLIOTECARIO o ADMIN
	\textbf{Postcondicion:} Reporte generado y entregado al solicitante
	
	\hypertarget{escenario-principal-de-exito-2}{%
		\subsection{Escenario principal de exito}\label{escenario-principal-de-exito-2}}
	
	\begin{enumerate}
		\def\labelenumi{\arabic{enumi}.}
		\tightlist
		\item
		El usuario solicita un tipo de reporte (prestamos diarios, libros mas solicitados, multas cobradas)
		\item
		El sistema ejecuta el stored procedure correspondiente
		\item
		El sistema recopila los datos agregados
		\item
		El sistema retorna los datos en formato JSON
	\end{enumerate}
	
	\hypertarget{extensiones-4}{%
		\subsection{Extensiones}\label{extensiones-4}}
	
	\hypertarget{a.-sin-datos-para-el-periodo}{%
		\subsubsection{2a. Sin datos para el periodo}\label{a.-sin-datos-para-el-periodo}}
	
	2a1. El sistema retorna un reporte vacio con total = 0
	
	\hypertarget{a.-error-de-base-de-datos-1}{%
		\subsubsection{3a. Error de base de datos}\label{a.-error-de-base-de-datos-1}}
	
	3a1. El sistema retorna HTTP 500 con ProblemDetail
	
	\hypertarget{requisitos-asociados-4}{%
		\subsection{Requisitos asociados}\label{requisitos-asociados-4}}
	
	\begin{itemize}
		\tightlist
		\item
		REQ-F-009 (Reporte de prestamos diarios)
		\item
		REQ-NF-006 (Control de acceso por roles)
	\end{itemize}
	
	\hypertarget{historias-de-usuario}{%
		\chapter{3. Historias de Usuario}\label{historias-de-usuario}}
	
	\hypertarget{hu-01-autenticaciuxf3n-de-usuarios}{%
		\section{HU-01: Autenticación de usuarios}\label{hu-01-autenticaciuxf3n-de-usuarios}}
	
	\textbf{Formato:} Connextra
	\textbf{Rol:} Usuario del sistema
	\textbf{Objetivo:} Iniciar sesión con credenciales
	\textbf{Beneficio:} Acceder a las funcionalidades del sistema según mi rol
	
	\textbf{Criterios INVEST:} Independent, Negotiable, Valuable, Estimable, Small, Testable
	
	\hypertarget{criterios-de-aceptaciuxf3n-gherkin}{%
		\subsection{Criterios de aceptación (Gherkin)}\label{criterios-de-aceptaciuxf3n-gherkin}}
	
	\begin{lstlisting}
		Feature: Autenticación
		Scenario: Inicio de sesión exitoso
		Given un usuario registrado con email "admin@dev.com" y contraseña "admin123"
		When el usuario envía POST /api/auth/login con email y contraseña válidos
		Then el sistema responde con HTTP 200
		And la respuesta incluye un JWT en cookie HttpOnly
		And el JWT contiene los claims iss, sub, aud, exp, nbf, iat, jti
		
		Scenario: Inicio de sesión con credenciales inválidas
		Given un usuario registrado con email "admin@dev.com"
		When el usuario envía POST /api/auth/login con contraseña incorrecta
		Then el sistema responde con HTTP 401
		And la respuesta sigue el formato ProblemDetail RFC 7807
	\end{lstlisting}
	
	\hypertarget{hu-02-registro-de-usuarios}{%
		\section{HU-02: Registro de usuarios}\label{hu-02-registro-de-usuarios}}
	
	\textbf{Formato:} Connextra
	\textbf{Rol:} Visitante
	\textbf{Objetivo:} Registrarme como usuario del sistema
	\textbf{Beneficio:} Obtener acceso como estudiante a los servicios bibliotecarios
	
	\hypertarget{criterios-de-aceptaciuxf3n-gherkin-1}{%
		\subsection{Criterios de aceptación (Gherkin)}\label{criterios-de-aceptaciuxf3n-gherkin-1}}
	
	\begin{lstlisting}
		Feature: Registro
		Scenario: Registro exitoso
		Given un visitante con datos válidos (nombre, email, contraseña)
		When envía POST /api/auth/register
		Then el sistema responde con HTTP 201
		And se crea un usuario con rol ESTUDIANTE
		And se emite un JWT en cookie
		
		Scenario: Registro con email duplicado
		Given un usuario existente con email "estudiante@test.com"
		When otro visitante intenta registrar el mismo email
		Then el sistema responde con HTTP 400
		And el mensaje indica que el email ya está registrado
	\end{lstlisting}
	
	\hypertarget{hu-03-cierre-de-sesiuxf3n}{%
		\section{HU-03: Cierre de sesión}\label{hu-03-cierre-de-sesiuxf3n}}
	
	\textbf{Formato:} Connextra
	\textbf{Rol:} Usuario autenticado
	\textbf{Objetivo:} Cerrar sesión
	\textbf{Beneficio:} Invalidar mi token actual y proteger mi cuenta
	
	\hypertarget{criterios-de-aceptaciuxf3n-gherkin-2}{%
		\subsection{Criterios de aceptación (Gherkin)}\label{criterios-de-aceptaciuxf3n-gherkin-2}}
	
	\begin{lstlisting}
		Feature: Cierre de sesión
		Scenario: Logout exitoso
		Given un usuario autenticado con JWT válido
		When envía POST /api/auth/logout
		Then el sistema responde con HTTP 200
		And el JTI del token se agrega a la blacklist
		And la cookie se elimina
		
		Scenario: Acceso post-logout
		Given un usuario cuyo JWT fue invalidado
		When intenta acceder a /api/libros
		Then el sistema responde con HTTP 401
	\end{lstlisting}
	
	\hypertarget{hu-04-gestiuxf3n-de-usuarios}{%
		\section{HU-04: Gestión de usuarios}\label{hu-04-gestiuxf3n-de-usuarios}}
	
	\textbf{Formato:} Connextra
	\textbf{Rol:} Administrador
	\textbf{Objetivo:} Gestionar usuarios del sistema
	\textbf{Beneficio:} Controlar quién tiene acceso y con qué rol
	
	\hypertarget{criterios-de-aceptaciuxf3n-gherkin-3}{%
		\subsection{Criterios de aceptación (Gherkin)}\label{criterios-de-aceptaciuxf3n-gherkin-3}}
	
	\begin{lstlisting}
		Feature: CRUD de usuarios
		Scenario: Listar usuarios
		Given un administrador autenticado
		When envía GET /api/usuarios
		Then recibe lista paginada de usuarios
		
		Scenario: Crear usuario (admin)
		Given un administrador autenticado
		When envía POST /api/usuarios con datos válidos
		Then el sistema crea el usuario con el rol especificado
		
		Scenario: Usuario no admin no puede gestionar usuarios
		Given un estudiante autenticado
		When envía GET /api/usuarios
		Then el sistema responde con HTTP 403
	\end{lstlisting}
	
	\hypertarget{hu-05-gestiuxf3n-de-material-bibliogruxe1fico}{%
		\section{HU-05: Gestión de material bibliográfico}\label{hu-05-gestiuxf3n-de-material-bibliogruxe1fico}}
	
	\textbf{Formato:} Connextra
	\textbf{Rol:} Bibliotecario/Administrador
	\textbf{Objetivo:} Gestionar libros, autores, editoriales y categorías
	\textbf{Beneficio:} Mantener actualizado el catálogo bibliotecario
	
	\hypertarget{criterios-de-aceptaciuxf3n-gherkin-4}{%
		\subsection{Criterios de aceptación (Gherkin)}\label{criterios-de-aceptaciuxf3n-gherkin-4}}
	
	\begin{lstlisting}
		Feature: CRUD de libros
		Scenario: Listar libros con paginación
		Given un usuario autenticado
		When envía GET /api/libros?page=0&size=10
		Then recibe página de libros activos
		
		Scenario: Buscar libros por título
		Given un usuario autenticado
		When envía GET /api/libros/buscar?q=programación
		Then recibe libros cuyo título contiene "programación"
		
		Scenario: Crear libro (bibliotecario)
		Given un bibliotecario autenticado
		When envía POST /api/libros con datos válidos
		Then el sistema crea el libro con ejemplares disponibles = totales
		
		Scenario: Estudiante no puede crear libros
		Given un estudiante autenticado
		When envía POST /api/libros
		Then el sistema responde con HTTP 403
	\end{lstlisting}
	
	\hypertarget{hu-06-prestamo-de-ejemplares}{%
		\section{HU-06: Prestamo de ejemplares}\label{hu-06-prestamo-de-ejemplares}}
	
	\textbf{Formato:} Connextra
	\textbf{Rol:} Bibliotecario
	\textbf{Objetivo:} Registrar el prestamo de un ejemplar a un usuario
	\textbf{Beneficio:} Controlar la circulacion del material bibliotecario
	
	\hypertarget{criterios-de-aceptacion-gherkin}{%
		\subsection{Criterios de aceptacion (Gherkin)}\label{criterios-de-aceptacion-gherkin}}
	
	\begin{lstlisting}
		Feature: Prestamo de ejemplares
		Scenario: Prestamo exitoso
		Given un usuario activo y un ejemplar disponible
		When el bibliotecario registra el prestamo
		Then el sistema crea el prestamo con fecha de vencimiento en 7 dias
		And el estado del ejemplar cambia a PRESTADO
		And los ejemplares disponibles del libro disminuyen en 1
		
		Scenario: Prestamo con ejemplar no disponible
		Given un ejemplar en estado PRESTADO
		When el bibliotecario intenta prestarlo
		Then el sistema rechaza la operacion con HTTP 400
		And el mensaje indica que el ejemplar no esta disponible
		
		Scenario: Prestamo con usuario que supera el limite
		Given un usuario con 5 prestamos activos
		When el bibliotecario intenta registrar un nuevo prestamo
		Then el sistema rechaza la operacion con HTTP 400
		And el mensaje indica limite de prestamos alcanzado
	\end{lstlisting}
	
	\hypertarget{hu-07-devolucion-de-prestamo}{%
		\section{HU-07: Devolucion de prestamo}\label{hu-07-devolucion-de-prestamo}}
	
	\textbf{Formato:} Connextra
	\textbf{Rol:} Bibliotecario
	\textbf{Objetivo:} Registrar la devolucion de un ejemplar prestado
	\textbf{Beneficio:} Actualizar el inventario y gestionar multas por retraso
	
	\hypertarget{criterios-de-aceptacion-gherkin-1}{%
		\subsection{Criterios de aceptacion (Gherkin)}\label{criterios-de-aceptacion-gherkin-1}}
	
	\begin{lstlisting}
		Feature: Devolucion de prestamo
		Scenario: Devolucion en fecha
		Given un prestamo activo cuya fecha de vencimiento no ha pasado
		When el bibliotecario registra la devolucion
		Then el sistema marca el prestamo como DEVUELTO
		And registra la fecha de devolucion
		And el ejemplar vuelve a estar DISPONIBLE
		And los ejemplares disponibles del libro aumentan en 1
		And no se genera ninguna multa
		
		Scenario: Devolucion con retraso
		Given un prestamo activo cuya fecha de vencimiento ya paso
		When el bibliotecario registra la devolucion
		Then el sistema genera una multa de $0.50 por dia de retraso
		And la multa queda registrada como PENDIENTE
	\end{lstlisting}
	
	\hypertarget{hu-08-renovacion-de-prestamo}{%
		\section{HU-08: Renovacion de prestamo}\label{hu-08-renovacion-de-prestamo}}
	
	\textbf{Formato:} Connextra
	\textbf{Rol:} Bibliotecario
	\textbf{Objetivo:} Renovar un prestamo activo
	\textbf{Beneficio:} Extender el periodo de prestamo sin devolver el material
	
	\hypertarget{criterios-de-aceptacion-gherkin-2}{%
		\subsection{Criterios de aceptacion (Gherkin)}\label{criterios-de-aceptacion-gherkin-2}}
	
	\begin{lstlisting}
		Feature: Renovacion de prestamo
		Scenario: Renovacion exitosa
		Given un prestamo activo
		When el bibliotecario solicita la renovacion
		Then el sistema marca el prestamo original como RENOVADO
		And crea un nuevo prestamo con fecha de vencimiento extendida 7 dias
		And el nuevo prestamo queda en estado ACTIVO
		
		Scenario: Renovacion de prestamo no activo
		Given un prestamo en estado DEVUELTO
		When el bibliotecario intenta renovarlo
		Then el sistema rechaza la operacion con HTTP 400
	\end{lstlisting}
	
	\hypertarget{hu-09-reportes-bibliotecarios}{%
		\section{HU-09: Reportes bibliotecarios}\label{hu-09-reportes-bibliotecarios}}
	
	\textbf{Formato:} Connextra
	\textbf{Rol:} Bibliotecario/Administrador
	\textbf{Objetivo:} Generar reportes de actividad
	\textbf{Beneficio:} Obtener datos para la toma de decisiones
	
	\hypertarget{criterios-de-aceptacion-gherkin-3}{%
		\subsection{Criterios de aceptacion (Gherkin)}\label{criterios-de-aceptacion-gherkin-3}}
	
	\begin{lstlisting}
		Feature: Reportes
		Scenario: Reporte de prestamos diarios
		Given un bibliotecario autenticado
		When solicita GET /api/reportes/prestamos-diarios
		Then recibe el total de prestamos del dia actual
		And el detalle de cada prestamo con usuario y libro
		
		Scenario: Reporte de libros mas solicitados
		Given un bibliotecario autenticado
		When solicita GET /api/reportes/libros-mas-solicitados
		Then recibe el top 10 de libros con mayor cantidad de prestamos
		
		Scenario: Reporte de multas cobradas
		Given un administrador autenticado
		When solicita GET /api/reportes/multas-cobradas
		Then recibe el total cobrado en el mes actual
	\end{lstlisting}
	
	\hypertarget{hu-10-dashboard-de-estadisticas}{%
		\section{HU-10: Dashboard de estadisticas}\label{hu-10-dashboard-de-estadisticas}}
	
	\textbf{Formato:} Connextra
	\textbf{Rol:} Usuario autenticado
	\textbf{Objetivo:} Visualizar estadisticas generales del sistema
	\textbf{Beneficio:} Obtener una vision rapida del estado de la biblioteca
	
	\hypertarget{criterios-de-aceptacion-gherkin-4}{%
		\subsection{Criterios de aceptacion (Gherkin)}\label{criterios-de-aceptacion-gherkin-4}}
	
	\begin{lstlisting}
		Feature: Dashboard
		Scenario: Visualizar estadisticas
		Given un usuario autenticado
		When solicita GET /api/dashboard/stats
		Then recibe libros prestados hoy
		And libros disponibles
		And estudiantes activos
		And reservas pendientes
		And multas pendientes
		And total acumulado de multas
	\end{lstlisting}
	
	\hypertarget{matriz-de-trazabilidad}{%
		\chapter{4. Matriz de Trazabilidad}\label{matriz-de-trazabilidad}}
	
	\begin{landscape}
		\footnotesize
		\setlength{\tabcolsep}{3pt}
		\renewcommand{\arraystretch}{1.3}
		\begin{longtable}{p{1.39cm}p{1.65cm}p{1.22cm}p{0.87cm}p{0.87cm}p{2.61cm}p{3.13cm}p{3.65cm}p{1.22cm}p{2.87cm}p{1.65cm}}
			\caption{Matriz de trazabilidad de requisitos}\\
			\toprule
			\textbf{ID} & \textbf{Tipo} & \textbf{Prioridad} & \textbf{HU} & \textbf{CU} & \textbf{Módulo} & \textbf{Endpoint API} & \textbf{Prueba automatizada} & \textbf{Acceso} & \textbf{Evidencia} & \textbf{Estado} \\
			\midrule
			\endfirsthead
			\toprule
			\textbf{ID} & \textbf{Tipo} & \textbf{Prioridad} & \textbf{HU} & \textbf{CU} & \textbf{Módulo} & \textbf{Endpoint API} & \textbf{Prueba automatizada} & \textbf{Acceso} & \textbf{Evidencia} & \textbf{Estado} \\
			\midrule
			\endhead
			\bottomrule
			\endfoot
			REQ-F-001 & Funcional & Must & HU-01 & CU-01 & \texttt{\footnotesize \seqsplit{AuthController}} & \texttt{\footnotesize \seqsplit{POST /api/auth/login}} & \texttt{\footnotesize \seqsplit{AuthControllerTest.login\_exitoso}} & CRUD-ORM & {\small --} & verificado \\
			REQ-F-001 & Funcional & Must & HU-01 & CU-01 & \texttt{\footnotesize \seqsplit{JwtTokenProvider}} & {\small --} & \texttt{\footnotesize \seqsplit{JwtTokenProviderTest.genera\_token\_con\_claims}} & CRUD-ORM & {\small --} & verificado \\
			REQ-F-002 & Funcional & Must & HU-02 & CU-01 & \texttt{\footnotesize \seqsplit{AuthController}} & \texttt{\footnotesize \seqsplit{POST /api/auth/register}} & \texttt{\footnotesize \seqsplit{AuthControllerTest.registro\_exitoso}} & CRUD-ORM & {\small --} & verificado \\
			REQ-F-003 & Funcional & Must & HU-03 & CU-01 & \texttt{\footnotesize \seqsplit{AuthController}} & \texttt{\footnotesize \seqsplit{POST /api/auth/logout}} & \texttt{\footnotesize \seqsplit{AuthControllerTest.logout\_invalida\_token}} & CRUD-ORM & {\small --} & verificado \\
			REQ-F-004 & Funcional & Must & HU-04 & {\small --} & \texttt{\footnotesize \seqsplit{UsuarioController}} & \texttt{\footnotesize \seqsplit{GET/POST/PUT/DELETE /api/usuarios}} & \texttt{\footnotesize \seqsplit{UsuarioControllerTest.crud}} & CRUD-ORM & {\small --} & verificado \\
			REQ-F-005 & Funcional & Must & HU-05 & CU-02 & \texttt{\footnotesize \seqsplit{LibroController}} & \texttt{\footnotesize \seqsplit{GET/POST/PUT/DELETE /api/libros}} & \texttt{\footnotesize \seqsplit{LibroControllerTest.crud}} & CRUD-ORM & {\small --} & verificado \\
			REQ-F-006 & Funcional & Must & HU-06 & CU-03 & \texttt{\footnotesize \seqsplit{PrestamoController}} & \texttt{\footnotesize \seqsplit{POST /api/prestamos}} & \texttt{\footnotesize \seqsplit{PrestamoServiceTest.crear\_prestamo}} & SP & \texttt{\footnotesize \seqsplit{docs/mediciones/perf/REPORT.md}} & verificado \\
			REQ-F-007 & Funcional & Must & HU-07 & CU-04 & \texttt{\footnotesize \seqsplit{PrestamoController}} & \texttt{\footnotesize \seqsplit{PUT /api/prestamos/\{id\}}/devolver} & \texttt{\footnotesize \seqsplit{PrestamoServiceTest.devolver\_prestamo}} & SP & {\small --} & verificado \\
			REQ-F-008 & Funcional & Should & HU-08 & {\small --} & \texttt{\footnotesize \seqsplit{PrestamoController}} & \texttt{\footnotesize \seqsplit{PUT /api/prestamos/\{id\}}/renovar} & {\small --} & SP & {\small --} & implementado \\
			REQ-F-009 & Funcional & Should & HU-09 & CU-05 & \texttt{\footnotesize \seqsplit{ReporteController}} & \texttt{\footnotesize \seqsplit{GET /api/reportes/prestamos-diarios}} & \texttt{\footnotesize \seqsplit{ReporteControllerTest.prestamosDiariosComoBibliotecario}} & SP & {\small --} & implementado \\
			REQ-F-010 & Funcional & Must & HU-10 & {\small --} & \texttt{\footnotesize \seqsplit{DashboardController}} & \texttt{\footnotesize \seqsplit{GET /api/dashboard/stats}} & {\small --} & SP & {\small --} & verificado \\
			REQ-F-011 & Funcional & Should & {\small --} & {\small --} & \texttt{\footnotesize \seqsplit{ReservaController}} & \texttt{\footnotesize \seqsplit{POST/DELETE /api/reservas/}} & {\small --} & CRUD-ORM & {\small --} & implementado \\
			REQ-NF-001 & No funcional & Must & HU-05 & {\small --} & \texttt{\footnotesize \seqsplit{LibroService}} & \texttt{\footnotesize \seqsplit{GET /api/libros}} & \texttt{\footnotesize \seqsplit{LibroPerfTest.p95\_200ms}} & CRUD-ORM & \texttt{\footnotesize \seqsplit{docs/mediciones/perf/REPORT.md}} & verificado \\
			REQ-NF-002 & No funcional & Must & HU-01 & CU-01 & \texttt{\footnotesize \seqsplit{SecurityConfig}} & {\small --} & \texttt{\footnotesize \seqsplit{SecurityTest.endpoints\_protegidos}} & CRUD-ORM & {\small --} & verificado \\
			REQ-NF-003 & No funcional & Must & {\small --} & {\small --} & \texttt{\footnotesize \seqsplit{GlobalExceptionHandler}} & {\small --} & \texttt{\footnotesize \seqsplit{SecurityTest.sql\_injection\_prevenido}} & CRUD-ORM & {\small --} & verificado \\
			REQ-NF-004 & No funcional & Must & {\small --} & {\small --} & {\small --} & \texttt{\footnotesize \seqsplit{docker-compose.yml}} & {\small --} & {\small --} & \texttt{\footnotesize \seqsplit{docs/mediciones/perf/startup-time.md}} & pendiente \\
			REQ-NF-005 & No funcional & Should & {\small --} & {\small --} & {\small --} & \texttt{\footnotesize \seqsplit{pom.xml}} & {\small --} & JaCoCo & {\small --} & verificado \\
			REQ-NF-006 & No funcional & Must & HU-04 & CU-01 & \texttt{\footnotesize \seqsplit{SecurityConfig}} & \texttt{\footnotesize \seqsplit{@PreAuthorize}} & \texttt{\footnotesize \seqsplit{ReporteControllerTest.prestamosDiariosComoEstudiante}} & CRUD-ORM & {\small --} & verificado \\
			REQ-NF-007 & No funcional & Must & HU-03 & {\small --} & \texttt{\footnotesize \seqsplit{JwtBlacklistService}} & {\small --} & \texttt{\footnotesize \seqsplit{JwtBlacklistServiceTest}} & CRUD-ORM & {\small --} & verificado \\
			REQ-R-001 & Rendimiento & Must & {\small --} & {\small --} & \texttt{\footnotesize \seqsplit{AuthController}} & \texttt{\footnotesize \seqsplit{POST /api/auth/login}} & {\small --} & CRUD-ORM & \texttt{\footnotesize \seqsplit{docs/mediciones/perf/}} & pendiente \\
			REQ-R-002 & Rendimiento & Must & {\small --} & {\small --} & \texttt{\footnotesize \seqsplit{LibroController}} & \texttt{\footnotesize \seqsplit{GET/POST/PUT/DELETE /api/libros}} & {\small --} & CRUD-ORM & \texttt{\footnotesize \seqsplit{docs/mediciones/perf/}} & pendiente \\
			REQ-R-003 & Rendimiento & Should & {\small --} & {\small --} & \texttt{\footnotesize \seqsplit{PrestamoController}} & \texttt{\footnotesize \seqsplit{POST /api/prestamos}} & {\small --} & SP & \texttt{\footnotesize \seqsplit{docs/mediciones/perf/}} & pendiente \\
			REQ-R-004 & Rendimiento & Should & {\small --} & {\small --} & \texttt{\footnotesize \seqsplit{LibroService}} & \texttt{\footnotesize \seqsplit{GET /api/libros}} & {\small --} & CRUD-ORM & \texttt{\footnotesize \seqsplit{docs/mediciones/perf/REPORT.md}} & pendiente \\
			\bottomrule
		\end{longtable}
	\end{landscape}
	
	\hypertarget{codigo-fuente}{%
		\chapter{5. Codigo Fuente}\label{codigo-fuente}}
	
	\hypertarget{jwttokenprovider.java}{%
		\section{JwtTokenProvider.java}\label{jwttokenprovider.java}}
	
	\passthrough{\lstinline!sigcb-qr-api/src/main/java/com/sigcbqr/security/JwtTokenProvider.java!}
	
	\begin{lstlisting}[language=Java]
		package com.sigcbqr.security;
		
		import io.jsonwebtoken.*;
		import io.jsonwebtoken.io.Decoders;
		import io.jsonwebtoken.security.Keys;
		import jakarta.servlet.http.Cookie;
		import jakarta.servlet.http.HttpServletRequest;
		import org.springframework.beans.factory.annotation.Value;
		import org.springframework.stereotype.Component;
		
		import javax.crypto.SecretKey;
		import java.util.Date;
		import java.util.UUID;
		
		@Component
		public class JwtTokenProvider {
			
			private final SecretKey secretKey;
			private final long expirationMs;
			private final String issuer;
			private final String audience;
			private final boolean secureCookie;
			private final JwtBlacklistService blacklistService;
			
			public JwtTokenProvider(
			@Value("${app.jwt.secret}") String secret,
			@Value("${app.jwt.expiration-ms}") long expirationMs,
			@Value("${app.jwt.issuer}") String issuer,
			@Value("${app.jwt.audience}") String audience,
			@Value("${app.jwt.secure-cookie}") boolean secureCookie,
			JwtBlacklistService blacklistService) {
				this.secretKey = Keys.hmacShaKeyFor(Decoders.BASE64.decode(secret));
				this.expirationMs = expirationMs;
				this.issuer = issuer;
				this.audience = audience;
				this.secureCookie = secureCookie;
				this.blacklistService = blacklistService;
			}
			
			public String generateToken(Long userId, String email, String rol) {
				Date now = new Date();
				Date expiryDate = new Date(now.getTime() + expirationMs);
				
				return Jwts.builder()
				.id(UUID.randomUUID().toString())
				.issuer(issuer)
				.subject(userId.toString())
				.audience().add(audience).and()
				.claim("email", email)
				.claim("rol", rol)
				.issuedAt(now)
				.notBefore(now)
				.expiration(expiryDate)
				.signWith(secretKey)
				.compact();
			}
			
			public Long getUserIdFromToken(String token) {
				Claims claims = parseToken(token);
				return Long.parseLong(claims.getSubject());
			}
			
			public String getEmailFromToken(String token) {
				Claims claims = parseToken(token);
				return claims.get("email", String.class);
			}
			
			public String getRolFromToken(String token) {
				Claims claims = parseToken(token);
				return claims.get("rol", String.class);
			}
			
			public String getJtiFromToken(String token) {
				Claims claims = parseToken(token);
				return claims.getId();
			}
			
			public boolean validateToken(String token) {
				try {
					Claims claims = parseToken(token);
					if (blacklistService.isBlacklisted(claims.getId())) {
						return false;
					}
					return true;
				} catch (JwtException | IllegalArgumentException e) {
					return false;
				}
			}
			
			private Claims parseToken(String token) {
				return Jwts.parser()
				.verifyWith(secretKey)
				.build()
				.parseSignedClaims(token)
				.getPayload();
			}
			
			public String extractTokenFromCookie(HttpServletRequest request) {
				Cookie[] cookies = request.getCookies();
				if (cookies != null) {
					for (Cookie cookie : cookies) {
						if ("access_token".equals(cookie.getName())) {
							return cookie.getValue();
						}
					}
				}
				return null;
			}
			
			public Cookie createAccessTokenCookie(String token) {
				Cookie cookie = new Cookie("access_token", token);
				cookie.setHttpOnly(true);
				cookie.setSecure(secureCookie);
				cookie.setPath("/");
				cookie.setMaxAge((int) (expirationMs / 1000));
				cookie.setAttribute("SameSite", "Strict");
				return cookie;
			}
			
			public Date getExpirationFromToken(String token) {
				return parseToken(token).getExpiration();
			}
			
			public SecretKey getSecretKey() {
				return secretKey;
			}
			
			public JwtBlacklistService getJwtBlacklistService() {
				return blacklistService;
			}
			
			public Cookie createLogoutCookie() {
				Cookie cookie = new Cookie("access_token", null);
				cookie.setHttpOnly(true);
				cookie.setSecure(secureCookie);
				cookie.setPath("/");
				cookie.setMaxAge(0);
				cookie.setAttribute("SameSite", "Strict");
				return cookie;
			}
		}
	\end{lstlisting}
	
	\hypertarget{jwtauthenticationfilter.java}{%
		\section{JwtAuthenticationFilter.java}\label{jwtauthenticationfilter.java}}
	
	\passthrough{\lstinline!sigcb-qr-api/src/main/java/com/sigcbqr/security/JwtAuthenticationFilter.java!}
	
	\begin{lstlisting}[language=Java]
		package com.sigcbqr.security;
		
		import jakarta.servlet.FilterChain;
		import jakarta.servlet.ServletException;
		import jakarta.servlet.http.HttpServletRequest;
		import jakarta.servlet.http.HttpServletResponse;
		import org.springframework.security.authentication.UsernamePasswordAuthenticationToken;
		import org.springframework.security.core.authority.SimpleGrantedAuthority;
		import org.springframework.security.core.context.SecurityContextHolder;
		import org.springframework.security.web.authentication.WebAuthenticationDetailsSource;
		import org.springframework.stereotype.Component;
		import org.springframework.util.StringUtils;
		import org.springframework.web.filter.OncePerRequestFilter;
		
		import java.io.IOException;
		import java.util.List;
		
		@Component
		public class JwtAuthenticationFilter extends OncePerRequestFilter {
			
			private final JwtTokenProvider tokenProvider;
			
			public JwtAuthenticationFilter(JwtTokenProvider tokenProvider) {
				this.tokenProvider = tokenProvider;
			}
			
			@Override
			protected void doFilterInternal(HttpServletRequest request,
			HttpServletResponse response,
			FilterChain filterChain) throws ServletException, IOException {
				
				String token = tokenProvider.extractTokenFromCookie(request);
				
				if (StringUtils.hasText(token) && tokenProvider.validateToken(token)) {
					Long userId = tokenProvider.getUserIdFromToken(token);
					String email = tokenProvider.getEmailFromToken(token);
					String rol = tokenProvider.getRolFromToken(token);
					
					List<SimpleGrantedAuthority> authorities = List.of(
					new SimpleGrantedAuthority("ROLE_" + rol)
					);
					
					UserPrincipal userPrincipal = new UserPrincipal(userId, email, null, rol, authorities);
					
					UsernamePasswordAuthenticationToken authentication =
					new UsernamePasswordAuthenticationToken(userPrincipal, null, authorities);
					authentication.setDetails(new WebAuthenticationDetailsSource().buildDetails(request));
					
					SecurityContextHolder.getContext().setAuthentication(authentication);
				}
				
				filterChain.doFilter(request, response);
			}
		}
	\end{lstlisting}
	
	\hypertarget{securityconfig.java}{%
		\section{SecurityConfig.java}\label{securityconfig.java}}
	
	\passthrough{\lstinline!sigcb-qr-api/src/main/java/com/sigcbqr/config/SecurityConfig.java!}
	
	\begin{lstlisting}[language=Java]
		package com.sigcbqr.config;
		
		import com.sigcbqr.security.JwtAuthenticationEntryPoint;
		import com.sigcbqr.security.JwtAuthenticationFilter;
		import org.springframework.context.annotation.Bean;
		import org.springframework.context.annotation.Configuration;
		import org.springframework.security.authentication.AuthenticationManager;
		import org.springframework.security.config.annotation.authentication.configuration.AuthenticationConfiguration;
		import org.springframework.security.config.annotation.method.configuration.EnableMethodSecurity;
		import org.springframework.security.config.annotation.web.builders.HttpSecurity;
		import org.springframework.security.config.annotation.web.configuration.EnableWebSecurity;
		import org.springframework.security.config.http.SessionCreationPolicy;
		import org.springframework.security.crypto.bcrypt.BCryptPasswordEncoder;
		import org.springframework.security.crypto.password.PasswordEncoder;
		import org.springframework.security.web.SecurityFilterChain;
		import org.springframework.security.web.authentication.UsernamePasswordAuthenticationFilter;
		
		@Configuration
		@EnableWebSecurity
		@EnableMethodSecurity
		public class SecurityConfig {
			
			private final JwtAuthenticationFilter jwtAuthenticationFilter;
			private final JwtAuthenticationEntryPoint jwtAuthenticationEntryPoint;
			
			public SecurityConfig(JwtAuthenticationFilter jwtAuthenticationFilter,
			JwtAuthenticationEntryPoint jwtAuthenticationEntryPoint) {
				this.jwtAuthenticationFilter = jwtAuthenticationFilter;
				this.jwtAuthenticationEntryPoint = jwtAuthenticationEntryPoint;
			}
			
			@Bean
			public SecurityFilterChain filterChain(HttpSecurity http) throws Exception {
				http
				.csrf(csrf -> csrf.disable())
				.cors(cors -> {})
				.exceptionHandling(ex -> ex.authenticationEntryPoint(jwtAuthenticationEntryPoint))
				.sessionManagement(session -> session.sessionCreationPolicy(SessionCreationPolicy.STATELESS))
				.authorizeHttpRequests(auth -> auth
				.requestMatchers("/api/auth/**").permitAll()
				.requestMatchers("/swagger-ui/**", "/api-docs/**", "/v3/api-docs/**").permitAll()
				.requestMatchers("/api/dashboard/**").authenticated()
				.requestMatchers("/api/**").authenticated()
				.anyRequest().permitAll()
				)
				.addFilterBefore(jwtAuthenticationFilter, UsernamePasswordAuthenticationFilter.class);
				
				return http.build();
			}
			
			@Bean
			public PasswordEncoder passwordEncoder() {
				return new BCryptPasswordEncoder();
			}
			
			@Bean
			public AuthenticationManager authenticationManager(AuthenticationConfiguration config) throws Exception {
				return config.getAuthenticationManager();
			}
		}
	\end{lstlisting}
	
	\hypertarget{jwtblacklistservice.java}{%
		\section{JwtBlacklistService.java}\label{jwtblacklistservice.java}}
	
	\passthrough{\lstinline!sigcb-qr-api/src/main/java/com/sigcbqr/security/JwtBlacklistService.java!}
	
	\begin{lstlisting}[language=Java]
		package com.sigcbqr.security;
		
		import com.sigcbqr.model.entity.JwtBlacklist;
		import com.sigcbqr.repository.JwtBlacklistRepository;
		import org.springframework.scheduling.annotation.Scheduled;
		import org.springframework.stereotype.Service;
		import org.springframework.transaction.annotation.Transactional;
		
		import java.time.LocalDateTime;
		
		@Service
		public class JwtBlacklistService {
			
			private final JwtBlacklistRepository repository;
			
			public JwtBlacklistService(JwtBlacklistRepository repository) {
				this.repository = repository;
			}
			
			@Transactional
			public void blacklist(String jti, LocalDateTime expiration) {
				if (!repository.existsByJti(jti)) {
					JwtBlacklist entry = JwtBlacklist.builder()
					.jti(jti)
					.fechaExpiracion(expiration)
					.build();
					repository.save(entry);
				}
			}
			
			public boolean isBlacklisted(String jti) {
				return repository.existsByJti(jti);
			}
			
			@Scheduled(cron = "0 0 */6 * * *")
			@Transactional
			public void cleanExpiredEntries() {
				repository.deleteByFechaExpiracionBefore(LocalDateTime.now());
			}
		}
	\end{lstlisting}
	
	\hypertarget{authcontroller.java}{%
		\section{AuthController.java}\label{authcontroller.java}}
	
	\passthrough{\lstinline!sigcb-qr-api/src/main/java/com/sigcbqr/controller/AuthController.java!}
	
	\begin{lstlisting}[language=Java]
		package com.sigcbqr.controller;
		
		import com.sigcbqr.model.dto.request.LoginRequest;
		import com.sigcbqr.model.dto.request.RegisterRequest;
		import com.sigcbqr.model.dto.response.ApiResponse;
		import com.sigcbqr.model.dto.response.LoginResponse;
		import com.sigcbqr.model.dto.response.UsuarioResponse;
		import com.sigcbqr.security.JwtTokenProvider;
		import com.sigcbqr.security.UserPrincipal;
		import com.sigcbqr.service.AuthService;
		import io.swagger.v3.oas.annotations.Operation;
		import io.swagger.v3.oas.annotations.tags.Tag;
		import jakarta.servlet.http.HttpServletRequest;
		import jakarta.servlet.http.HttpServletResponse;
		import jakarta.validation.Valid;
		import org.springframework.http.ResponseEntity;
		import org.springframework.security.core.annotation.AuthenticationPrincipal;
		import org.springframework.web.bind.annotation.*;
		
		import java.time.LocalDateTime;
		import java.util.Date;
		
		@RestController
		@RequestMapping("/api/auth")
		@Tag(name = "Autenticación", description = "Endpoints de autenticación con JWT en cookie HttpOnly")
		public class AuthController {
			
			private final AuthService authService;
			private final JwtTokenProvider tokenProvider;
			
			public AuthController(AuthService authService, JwtTokenProvider tokenProvider) {
				this.authService = authService;
				this.tokenProvider = tokenProvider;
			}
			
			@PostMapping("/login")
			@Operation(summary = "Iniciar sesión", description = "Autentica al usuario y devuelve un JWT en cookie HttpOnly")
			public ResponseEntity<ApiResponse> login(@Valid @RequestBody LoginRequest request,
			HttpServletResponse response) {
				LoginResponse loginResponse = authService.login(request);
				
				String token = tokenProvider.generateToken(
				loginResponse.getId(),
				loginResponse.getEmail(),
				loginResponse.getRol()
				);
				
				var cookie = tokenProvider.createAccessTokenCookie(token);
				response.addCookie(cookie);
				
				return ResponseEntity.ok(ApiResponse.success("Inicio de sesión exitoso", loginResponse));
			}
			
			@PostMapping("/register")
			@Operation(summary = "Registrar usuario", description = "Registra un nuevo usuario y devuelve JWT en cookie")
			public ResponseEntity<ApiResponse> register(@Valid @RequestBody RegisterRequest request,
			HttpServletResponse response) {
				LoginResponse loginResponse = authService.register(request);
				
				String token = tokenProvider.generateToken(
				loginResponse.getId(),
				loginResponse.getEmail(),
				loginResponse.getRol()
				);
				
				var cookie = tokenProvider.createAccessTokenCookie(token);
				response.addCookie(cookie);
				
				return ResponseEntity.ok(ApiResponse.success("Registro exitoso", loginResponse));
			}
			
			@PostMapping("/logout")
			@Operation(summary = "Cerrar sesión", description = "Invalida el JWT actual agregando su JTI a la blacklist")
			public ResponseEntity<ApiResponse> logout(HttpServletRequest request,
			HttpServletResponse response) {
				String token = tokenProvider.extractTokenFromCookie(request);
				if (token != null && tokenProvider.validateToken(token)) {
					Date expiration = tokenProvider.getExpirationFromToken(token);
					authService.logout(token, new java.sql.Timestamp(expiration.getTime()).toLocalDateTime());
				}
				
				var cookie = tokenProvider.createLogoutCookie();
				response.addCookie(cookie);
				
				return ResponseEntity.ok(ApiResponse.success("Sesión cerrada exitosamente", null));
			}
			
			@GetMapping("/me")
			@Operation(summary = "Usuario actual", description = "Obtiene los datos del usuario autenticado")
			public ResponseEntity<ApiResponse> me(@AuthenticationPrincipal UserPrincipal userPrincipal) {
				var usuario = authService.getCurrentUser(userPrincipal.id());
				var response = UsuarioResponse.builder()
				.id(usuario.getId())
				.nombre(usuario.getNombre())
				.email(usuario.getEmail())
				.telefono(usuario.getTelefono())
				.rol(usuario.getRol().getNombre())
				.activo(usuario.getActivo())
				.createdAt(usuario.getCreatedAt())
				.build();
				return ResponseEntity.ok(ApiResponse.success("Usuario actual", response));
			}
		}
	\end{lstlisting}
	
	\hypertarget{authservice.java}{%
		\section{AuthService.java}\label{authservice.java}}
	
	\passthrough{\lstinline!sigcb-qr-api/src/main/java/com/sigcbqr/service/AuthService.java!}
	
	\begin{lstlisting}[language=Java]
		package com.sigcbqr.service;
		
		import com.sigcbqr.exception.BadRequestException;
		import com.sigcbqr.exception.ResourceNotFoundException;
		import com.sigcbqr.model.dto.request.LoginRequest;
		import com.sigcbqr.model.dto.request.RegisterRequest;
		import com.sigcbqr.model.dto.response.LoginResponse;
		import com.sigcbqr.model.entity.Rol;
		import com.sigcbqr.model.entity.Usuario;
		import com.sigcbqr.repository.RolRepository;
		import com.sigcbqr.repository.UsuarioRepository;
		import com.sigcbqr.security.JwtTokenProvider;
		import org.springframework.security.authentication.AuthenticationManager;
		import org.springframework.security.authentication.UsernamePasswordAuthenticationToken;
		import org.springframework.security.crypto.password.PasswordEncoder;
		import org.springframework.stereotype.Service;
		import org.springframework.transaction.annotation.Transactional;
		
		import java.time.LocalDateTime;
		
		@Service
		public class AuthService {
			
			private final AuthenticationManager authenticationManager;
			private final JwtTokenProvider tokenProvider;
			private final UsuarioRepository usuarioRepository;
			private final RolRepository rolRepository;
			private final PasswordEncoder passwordEncoder;
			
			public AuthService(AuthenticationManager authenticationManager,
			JwtTokenProvider tokenProvider,
			UsuarioRepository usuarioRepository,
			RolRepository rolRepository,
			PasswordEncoder passwordEncoder) {
				this.authenticationManager = authenticationManager;
				this.tokenProvider = tokenProvider;
				this.usuarioRepository = usuarioRepository;
				this.rolRepository = rolRepository;
				this.passwordEncoder = passwordEncoder;
			}
			
			public LoginResponse login(LoginRequest request) {
				authenticationManager.authenticate(
				new UsernamePasswordAuthenticationToken(request.getEmail(), request.getPassword())
				);
				
				Usuario usuario = usuarioRepository.findByEmail(request.getEmail())
				.orElseThrow(() -> new ResourceNotFoundException("Usuario no encontrado"));
				
				String token = tokenProvider.generateToken(usuario.getId(), usuario.getEmail(),
				usuario.getRol().getNombre());
				
				return LoginResponse.builder()
				.id(usuario.getId())
				.nombre(usuario.getNombre())
				.email(usuario.getEmail())
				.rol(usuario.getRol().getNombre())
				.mensaje("Inicio de sesión exitoso")
				.build();
			}
			
			@Transactional
			public LoginResponse register(RegisterRequest request) {
				if (usuarioRepository.existsByEmail(request.getEmail())) {
					throw new BadRequestException("El correo ya está registrado");
				}
				
				Rol rol;
				if (request.getRolId() != null) {
					rol = rolRepository.findById(request.getRolId())
					.orElseThrow(() -> new ResourceNotFoundException("Rol no encontrado"));
				} else {
					rol = rolRepository.findByNombre("ESTUDIANTE")
					.orElseThrow(() -> new ResourceNotFoundException("Rol ESTUDIANTE no encontrado"));
				}
				
				Usuario usuario = Usuario.builder()
				.nombre(request.getNombre())
				.email(request.getEmail())
				.password(passwordEncoder.encode(request.getPassword()))
				.telefono(request.getTelefono())
				.activo(true)
				.rol(rol)
				.build();
				
				usuario = usuarioRepository.save(usuario);
				
				String token = tokenProvider.generateToken(usuario.getId(), usuario.getEmail(),
				usuario.getRol().getNombre());
				
				return LoginResponse.builder()
				.id(usuario.getId())
				.nombre(usuario.getNombre())
				.email(usuario.getEmail())
				.rol(usuario.getRol().getNombre())
				.mensaje("Registro exitoso")
				.build();
			}
			
			public void logout(String token, LocalDateTime expiration) {
				String jti = tokenProvider.getJtiFromToken(token);
				tokenProvider.getJwtBlacklistService().blacklist(jti, expiration);
			}
			
			public Usuario getCurrentUser(Long userId) {
				return usuarioRepository.findById(userId)
				.orElseThrow(() -> new ResourceNotFoundException("Usuario", userId));
			}
		}
	\end{lstlisting}
	
	\hypertarget{prestamocontroller.java}{%
		\section{PrestamoController.java}\label{prestamocontroller.java}}
	
	\passthrough{\lstinline!sigcb-qr-api/src/main/java/com/sigcbqr/controller/PrestamoController.java!}
	
	\begin{lstlisting}[language=Java]
		package com.sigcbqr.controller;
		
		import com.sigcbqr.model.dto.request.PrestamoRequest;
		import com.sigcbqr.model.dto.response.ApiResponse;
		import com.sigcbqr.model.dto.response.PageResponse;
		import com.sigcbqr.model.dto.response.PrestamoResponse;
		import com.sigcbqr.service.PrestamoService;
		import io.swagger.v3.oas.annotations.Operation;
		import io.swagger.v3.oas.annotations.tags.Tag;
		import jakarta.validation.Valid;
		import org.springframework.data.domain.Pageable;
		import org.springframework.data.web.PageableDefault;
		import org.springframework.http.ResponseEntity;
		import org.springframework.security.access.prepost.PreAuthorize;
		import org.springframework.web.bind.annotation.*;
		
		@RestController
		@RequestMapping("/api/prestamos")
		@Tag(name = "Préstamos", description = "Gestión de préstamos, devoluciones y renovaciones")
		public class PrestamoController {
			
			private final PrestamoService prestamoService;
			
			public PrestamoController(PrestamoService prestamoService) {
				this.prestamoService = prestamoService;
			}
			
			@GetMapping
			@Operation(summary = "Listar préstamos", description = "Obtiene todos los préstamos con paginación")
			@PreAuthorize("hasAnyRole('ADMIN', 'BIBLIOTECARIO')")
			public ResponseEntity<PageResponse<PrestamoResponse>> listar(
			@PageableDefault(size = 10, sort = "fechaPrestamo") Pageable pageable) {
				var page = prestamoService.listar(pageable);
				return ResponseEntity.ok(PageResponse.from(page));
			}
			
			@GetMapping("/usuario/{usuarioId}")
			@Operation(summary = "Préstamos por usuario", description = "Obtiene los préstamos de un usuario específico")
			public ResponseEntity<PageResponse<PrestamoResponse>> listarPorUsuario(
			@PathVariable Long usuarioId,
			@PageableDefault(size = 10) Pageable pageable) {
				var page = prestamoService.listarPorUsuario(usuarioId, pageable);
				return ResponseEntity.ok(PageResponse.from(page));
			}
			
			@GetMapping("/{id}")
			@Operation(summary = "Obtener préstamo", description = "Obtiene un préstamo por su ID")
			public ResponseEntity<ApiResponse> obtener(@PathVariable Long id) {
				var prestamo = prestamoService.obtener(id);
				return ResponseEntity.ok(ApiResponse.success("Préstamo encontrado", prestamo));
			}
			
			@PostMapping
			@PreAuthorize("hasAnyRole('ADMIN', 'BIBLIOTECARIO')")
			@Operation(summary = "Crear préstamo", description = "Registra un nuevo préstamo")
			public ResponseEntity<ApiResponse> crear(@Valid @RequestBody PrestamoRequest request) {
				var prestamo = prestamoService.crear(request);
				return ResponseEntity.ok(ApiResponse.created("Préstamo registrado", prestamo));
			}
			
			@PutMapping("/{id}/devolver")
			@PreAuthorize("hasAnyRole('ADMIN', 'BIBLIOTECARIO')")
			@Operation(summary = "Devolver libro", description = "Registra la devolución de un préstamo")
			public ResponseEntity<ApiResponse> devolver(@PathVariable Long id) {
				var prestamo = prestamoService.devolver(id);
				return ResponseEntity.ok(ApiResponse.success("Devolución registrada", prestamo));
			}
			
			@PutMapping("/{id}/renovar")
			@PreAuthorize("hasAnyRole('ADMIN', 'BIBLIOTECARIO')")
			@Operation(summary = "Renovar préstamo", description = "Renueva un préstamo activo")
			public ResponseEntity<ApiResponse> renovar(@PathVariable Long id) {
				var prestamo = prestamoService.renovar(id);
				return ResponseEntity.ok(ApiResponse.success("Préstamo renovado", prestamo));
			}
		}
	\end{lstlisting}
	
	\hypertarget{prestamoservice.java}{%
		\section{PrestamoService.java}\label{prestamoservice.java}}
	
	\passthrough{\lstinline!sigcb-qr-api/src/main/java/com/sigcbqr/service/PrestamoService.java!}
	
	\begin{lstlisting}[language=Java]
		package com.sigcbqr.service;
		
		import com.sigcbqr.exception.BadRequestException;
		import com.sigcbqr.exception.ResourceNotFoundException;
		import com.sigcbqr.model.dto.request.PrestamoRequest;
		import com.sigcbqr.model.dto.response.PrestamoResponse;
		import com.sigcbqr.model.entity.*;
		import com.sigcbqr.repository.*;
		import org.springframework.data.domain.Page;
		import org.springframework.data.domain.Pageable;
		import org.springframework.stereotype.Service;
		import org.springframework.transaction.annotation.Transactional;
		
		import java.time.LocalDateTime;
		
		@Service
		public class PrestamoService {
			
			private final PrestamoRepository prestamoRepository;
			private final UsuarioRepository usuarioRepository;
			private final InventarioRepository inventarioRepository;
			private final LibroRepository libroRepository;
			
			public PrestamoService(PrestamoRepository prestamoRepository,
			UsuarioRepository usuarioRepository,
			InventarioRepository inventarioRepository,
			LibroRepository libroRepository) {
				this.prestamoRepository = prestamoRepository;
				this.usuarioRepository = usuarioRepository;
				this.inventarioRepository = inventarioRepository;
				this.libroRepository = libroRepository;
			}
			
			public Page<PrestamoResponse> listar(Pageable pageable) {
				return prestamoRepository.findAll(pageable).map(this::toResponse);
			}
			
			public Page<PrestamoResponse> listarPorUsuario(Long usuarioId, Pageable pageable) {
				return prestamoRepository.findByUsuarioIdOrderByFechaPrestamoDesc(usuarioId, pageable)
				.map(this::toResponse);
			}
			
			public PrestamoResponse obtener(Long id) {
				Prestamo prestamo = prestamoRepository.findById(id)
				.orElseThrow(() -> new ResourceNotFoundException("Préstamo", id));
				return toResponse(prestamo);
			}
			
			@Transactional
			public PrestamoResponse crear(PrestamoRequest request) {
				Usuario usuario = usuarioRepository.findById(request.getUsuarioId())
				.orElseThrow(() -> new ResourceNotFoundException("Usuario", request.getUsuarioId()));
				
				Inventario inventario = inventarioRepository.findById(request.getInventarioId())
				.orElseThrow(() -> new ResourceNotFoundException("Ejemplar", request.getInventarioId()));
				
				if (!"DISPONIBLE".equals(inventario.getEstado())) {
					throw new BadRequestException("El ejemplar no está disponible");
				}
				
				long prestamosActivos = prestamoRepository.countByUsuarioIdAndEstado(request.getUsuarioId(), "ACTIVO");
				if (prestamosActivos >= 5) {
					throw new BadRequestException("El usuario tiene demasiados préstamos activos");
				}
				
				Prestamo prestamo = Prestamo.builder()
				.usuario(usuario)
				.inventario(inventario)
				.fechaPrestamo(LocalDateTime.now())
				.fechaVencimiento(LocalDateTime.now().plusDays(7))
				.estado("ACTIVO")
				.build();
				
				inventario.setEstado("PRESTADO");
				inventarioRepository.save(inventario);
				
				Libro libro = inventario.getLibro();
				libro.setEjemplaresDisponibles(libro.getEjemplaresDisponibles() - 1);
				libroRepository.save(libro);
				
				prestamo = prestamoRepository.save(prestamo);
				return toResponse(prestamo);
			}
			
			@Transactional
			public PrestamoResponse devolver(Long id) {
				Prestamo prestamo = prestamoRepository.findById(id)
				.orElseThrow(() -> new ResourceNotFoundException("Préstamo", id));
				
				if ("DEVUELTO".equals(prestamo.getEstado())) {
					throw new BadRequestException("El préstamo ya fue devuelto");
				}
				
				prestamo.setFechaDevolucion(LocalDateTime.now());
				prestamo.setEstado("DEVUELTO");
				
				Inventario inventario = prestamo.getInventario();
				inventario.setEstado("DISPONIBLE");
				inventarioRepository.save(inventario);
				
				Libro libro = inventario.getLibro();
				libro.setEjemplaresDisponibles(libro.getEjemplaresDisponibles() + 1);
				libroRepository.save(libro);
				
				// Generar multa si está vencido
				if (prestamo.getFechaVencimiento().isBefore(LocalDateTime.now())) {
					// Multa lógica here - handled by MultaService
				}
				
				prestamo = prestamoRepository.save(prestamo);
				return toResponse(prestamo);
			}
			
			@Transactional
			public PrestamoResponse renovar(Long id) {
				Prestamo prestamo = prestamoRepository.findById(id)
				.orElseThrow(() -> new ResourceNotFoundException("Préstamo", id));
				
				if (!"ACTIVO".equals(prestamo.getEstado())) {
					throw new BadRequestException("Solo se pueden renovar préstamos activos");
				}
				
				prestamo.setEstado("RENOVADO");
				prestamo.setObservaciones("Renovado - nueva fecha: " + LocalDateTime.now().plusDays(7));
				prestamoRepository.save(prestamo);
				
				Prestamo nuevoPrestamo = Prestamo.builder()
				.usuario(prestamo.getUsuario())
				.inventario(prestamo.getInventario())
				.fechaPrestamo(LocalDateTime.now())
				.fechaVencimiento(LocalDateTime.now().plusDays(7))
				.estado("ACTIVO")
				.observaciones("Renovación del préstamo #" + prestamo.getId())
				.build();
				
				nuevoPrestamo = prestamoRepository.save(nuevoPrestamo);
				return toResponse(nuevoPrestamo);
			}
			
			private PrestamoResponse toResponse(Prestamo prestamo) {
				return PrestamoResponse.builder()
				.id(prestamo.getId())
				.usuarioNombre(prestamo.getUsuario().getNombre())
				.libroTitulo(prestamo.getInventario().getLibro().getTitulo())
				.codigoEjemplar(prestamo.getInventario().getCodigoEjemplar())
				.fechaPrestamo(prestamo.getFechaPrestamo())
				.fechaVencimiento(prestamo.getFechaVencimiento())
				.fechaDevolucion(prestamo.getFechaDevolucion())
				.estado(prestamo.getEstado())
				.observaciones(prestamo.getObservaciones())
				.build();
			}
		}
	\end{lstlisting}
	
	\hypertarget{cacheconfig.java}{%
		\section{CacheConfig.java}\label{cacheconfig.java}}
	
	\passthrough{\lstinline!sigcb-qr-api/src/main/java/com/sigcbqr/config/CacheConfig.java!}
	
	\begin{lstlisting}[language=Java]
		package com.sigcbqr.config;
		
		import org.springframework.beans.factory.annotation.Value;
		import org.springframework.cache.CacheManager;
		import org.springframework.cache.annotation.EnableCaching;
		import org.springframework.context.annotation.Bean;
		import org.springframework.context.annotation.Configuration;
		import org.springframework.data.redis.cache.RedisCacheConfiguration;
		import org.springframework.data.redis.cache.RedisCacheManager;
		import org.springframework.data.redis.connection.RedisConnectionFactory;
		import org.springframework.data.redis.serializer.GenericJackson2JsonRedisSerializer;
		import org.springframework.data.redis.serializer.RedisSerializationContext;
		
		import java.time.Duration;
		
		@Configuration
		@EnableCaching
		public class CacheConfig {
			
			@Value("${app.cache.default-ttl-seconds}")
			private int defaultTtlSeconds;
			
			@Bean
			public CacheManager cacheManager(RedisConnectionFactory connectionFactory) {
				RedisCacheConfiguration config = RedisCacheConfiguration.defaultCacheConfig()
				.entryTtl(Duration.ofSeconds(defaultTtlSeconds))
				.disableCachingNullValues()
				.serializeValuesWith(
				RedisSerializationContext.SerializationPair.fromSerializer(
				new GenericJackson2JsonRedisSerializer()));
				
				return RedisCacheManager.builder(connectionFactory)
				.cacheDefaults(config)
				.build();
			}
		}
	\end{lstlisting}
	
	\hypertarget{globalexceptionhandler.java}{%
		\section{GlobalExceptionHandler.java}\label{globalexceptionhandler.java}}
	
	\passthrough{\lstinline!sigcb-qr-api/src/main/java/com/sigcbqr/exception/GlobalExceptionHandler.java!}
	
	\begin{lstlisting}[language=Java]
		package com.sigcbqr.exception;
		
		import org.springframework.http.HttpStatus;
		import org.springframework.http.ProblemDetail;
		import org.springframework.security.authentication.BadCredentialsException;
		import org.springframework.validation.FieldError;
		import org.springframework.web.bind.MethodArgumentNotValidException;
		import org.springframework.web.bind.annotation.ExceptionHandler;
		import org.springframework.web.bind.annotation.RestControllerAdvice;
		
		import java.net.URI;
		import java.util.HashMap;
		import java.util.Map;
		
		@RestControllerAdvice
		public class GlobalExceptionHandler {
			
			public static final String BASE_ERROR_URL = "https://api.sigcbqr.com/errors";
			
			@ExceptionHandler(ResourceNotFoundException.class)
			public ProblemDetail handleNotFound(ResourceNotFoundException ex) {
				ProblemDetail pd = ProblemDetail.forStatusAndDetail(HttpStatus.NOT_FOUND, ex.getMessage());
				pd.setType(URI.create(BASE_ERROR_URL + "/not-found"));
				pd.setTitle("Recurso no encontrado");
				pd.setProperty("timestamp", System.currentTimeMillis());
				return pd;
			}
			
			@ExceptionHandler(BadRequestException.class)
			public ProblemDetail handleBadRequest(BadRequestException ex) {
				ProblemDetail pd = ProblemDetail.forStatusAndDetail(HttpStatus.BAD_REQUEST, ex.getMessage());
				pd.setType(URI.create(BASE_ERROR_URL + "/bad-request"));
				pd.setTitle("Solicitud inválida");
				pd.setProperty("timestamp", System.currentTimeMillis());
				return pd;
			}
			
			@ExceptionHandler(BadCredentialsException.class)
			public ProblemDetail handleBadCredentials(BadCredentialsException ex) {
				ProblemDetail pd = ProblemDetail.forStatusAndDetail(HttpStatus.UNAUTHORIZED, "Credenciales inválidas");
				pd.setType(URI.create(BASE_ERROR_URL + "/unauthorized"));
				pd.setTitle("No autorizado");
				pd.setProperty("timestamp", System.currentTimeMillis());
				return pd;
			}
			
			@ExceptionHandler(MethodArgumentNotValidException.class)
			public ProblemDetail handleValidation(MethodArgumentNotValidException ex) {
				Map<String, String> errors = new HashMap<>();
				for (FieldError error : ex.getBindingResult().getFieldErrors()) {
					errors.put(error.getField(), error.getDefaultMessage());
				}
				ProblemDetail pd = ProblemDetail.forStatusAndDetail(HttpStatus.BAD_REQUEST, "Error de validación");
				pd.setType(URI.create(BASE_ERROR_URL + "/validation"));
				pd.setTitle("Error de validación");
				pd.setProperty("errors", errors);
				pd.setProperty("timestamp", System.currentTimeMillis());
				return pd;
			}
			
			@ExceptionHandler(Exception.class)
			public ProblemDetail handleGeneral(Exception ex) {
				ProblemDetail pd = ProblemDetail.forStatusAndDetail(HttpStatus.INTERNAL_SERVER_ERROR,
				"Error interno del servidor: " + ex.getMessage());
				pd.setType(URI.create(BASE_ERROR_URL + "/internal-error"));
				pd.setTitle("Error interno");
				pd.setProperty("timestamp", System.currentTimeMillis());
				return pd;
			}
		}
	\end{lstlisting}
	
	\hypertarget{prestamo.java}{%
		\section{Prestamo.java}\label{prestamo.java}}
	
	\passthrough{\lstinline!sigcb-qr-api/src/main/java/com/sigcbqr/model/entity/Prestamo.java!}
	
	\begin{lstlisting}[language=Java]
		package com.sigcbqr.model.entity;
		
		import com.fasterxml.jackson.annotation.JsonIgnoreProperties;
		import jakarta.persistence.*;
		import lombok.*;
		
		import java.time.LocalDateTime;
		
		@Entity
		@Table(name = "prestamos")
		@Getter @Setter @NoArgsConstructor @AllArgsConstructor @Builder
		@JsonIgnoreProperties({"hibernateLazyInitializer", "handler"})
		public class Prestamo {
			
			@Id
			@GeneratedValue(strategy = GenerationType.IDENTITY)
			private Long id;
			
			@ManyToOne(fetch = FetchType.LAZY)
			@JoinColumn(name = "usuario_id", nullable = false)
			private Usuario usuario;
			
			@ManyToOne(fetch = FetchType.LAZY)
			@JoinColumn(name = "inventario_id", nullable = false)
			private Inventario inventario;
			
			@Column(name = "fecha_prestamo")
			private LocalDateTime fechaPrestamo;
			
			@Column(name = "fecha_vencimiento", nullable = false)
			private LocalDateTime fechaVencimiento;
			
			@Column(name = "fecha_devolucion")
			private LocalDateTime fechaDevolucion;
			
			@Column(length = 30)
			private String estado;
			
			@Column(columnDefinition = "TEXT")
			private String observaciones;
			
			@Column(name = "created_at")
			private LocalDateTime createdAt;
			
			@Column(name = "updated_at")
			private LocalDateTime updatedAt;
			
			@PrePersist
			protected void onCreate() {
				createdAt = LocalDateTime.now();
				updatedAt = LocalDateTime.now();
				if (estado == null) estado = "ACTIVO";
				if (fechaPrestamo == null) fechaPrestamo = LocalDateTime.now();
			}
			
			@PreUpdate
			protected void onUpdate() {
				updatedAt = LocalDateTime.now();
			}
		}
	\end{lstlisting}
	
	\hypertarget{libro.java}{%
		\section{Libro.java}\label{libro.java}}
	
	\passthrough{\lstinline!sigcb-qr-api/src/main/java/com/sigcbqr/model/entity/Libro.java!}
	
	\begin{lstlisting}[language=Java]
		package com.sigcbqr.model.entity;
		
		import jakarta.persistence.*;
		import lombok.*;
		
		import java.time.LocalDateTime;
		import java.util.HashSet;
		import java.util.Set;
		
		@Entity
		@Table(name = "libros")
		@Getter @Setter @NoArgsConstructor @AllArgsConstructor @Builder
		public class Libro {
			
			@Id
			@GeneratedValue(strategy = GenerationType.IDENTITY)
			private Long id;
			
			@Column(nullable = false, length = 300)
			private String titulo;
			
			@Column(unique = true, length = 20)
			private String isbn;
			
			@Column(name = "anio_publicacion")
			private Integer anioPublicacion;
			
			@Column(length = 50)
			private String edicion;
			
			@Column(name = "ejemplares_totales")
			private Integer ejemplaresTotales;
			
			@Column(name = "ejemplares_disponibles")
			private Integer ejemplaresDisponibles;
			
			@Column(length = 100)
			private String ubicacion;
			
			@Column(columnDefinition = "TEXT")
			private String descripcion;
			
			private String portada;
			
			private Boolean activo;
			
			@ManyToOne(fetch = FetchType.LAZY)
			@JoinColumn(name = "categoria_id")
			private Categoria categoria;
			
			@ManyToOne(fetch = FetchType.LAZY)
			@JoinColumn(name = "editorial_id")
			private Editorial editorial;
			
			@ManyToMany(fetch = FetchType.LAZY)
			@JoinTable(
			name = "libro_autores",
			joinColumns = @JoinColumn(name = "libro_id"),
			inverseJoinColumns = @JoinColumn(name = "autor_id")
			)
			@Builder.Default
			private Set<Autor> autores = new HashSet<>();
			
			@Column(name = "created_at")
			private LocalDateTime createdAt;
			
			@Column(name = "updated_at")
			private LocalDateTime updatedAt;
			
			@PrePersist
			protected void onCreate() {
				createdAt = LocalDateTime.now();
				updatedAt = LocalDateTime.now();
				if (activo == null) activo = true;
				if (ejemplaresTotales == null) ejemplaresTotales = 1;
				if (ejemplaresDisponibles == null) ejemplaresDisponibles = ejemplaresTotales;
			}
			
			@PreUpdate
			protected void onUpdate() {
				updatedAt = LocalDateTime.now();
			}
		}
	\end{lstlisting}
	
	\hypertarget{usuario.java}{%
		\section{Usuario.java}\label{usuario.java}}
	
	\passthrough{\lstinline!sigcb-qr-api/src/main/java/com/sigcbqr/model/entity/Usuario.java!}
	
	\begin{lstlisting}[language=Java]
		package com.sigcbqr.model.entity;
		
		import jakarta.persistence.*;
		import lombok.*;
		
		import java.time.LocalDateTime;
		
		@Entity
		@Table(name = "usuarios")
		@Getter @Setter @NoArgsConstructor @AllArgsConstructor @Builder
		public class Usuario {
			
			@Id
			@GeneratedValue(strategy = GenerationType.IDENTITY)
			private Long id;
			
			@Column(nullable = false, length = 100)
			private String nombre;
			
			@Column(nullable = false, unique = true, length = 150)
			private String email;
			
			@Column(nullable = false)
			private String password;
			
			@Column(length = 20)
			private String telefono;
			
			private String foto;
			
			@Column(nullable = false)
			private Boolean activo;
			
			@ManyToOne(fetch = FetchType.EAGER)
			@JoinColumn(name = "rol_id", nullable = false)
			private Rol rol;
			
			@Column(name = "created_at")
			private LocalDateTime createdAt;
			
			@Column(name = "updated_at")
			private LocalDateTime updatedAt;
			
			@PrePersist
			protected void onCreate() {
				createdAt = LocalDateTime.now();
				updatedAt = LocalDateTime.now();
				if (activo == null) activo = true;
			}
			
			@PreUpdate
			protected void onUpdate() {
				updatedAt = LocalDateTime.now();
			}
		}
	\end{lstlisting}
	
	\hypertarget{sp_crear_prestamo.sql}{%
		\section{sp\_crear\_prestamo.sql}\label{sp_crear_prestamo.sql}}
	
	\passthrough{\lstinline!db/procs/sp\_crear\_prestamo.sql!}
	
	\begin{lstlisting}[language=SQL]
		CREATE OR REPLACE PROCEDURE sp_crear_prestamo(
		IN p_usuario_id BIGINT,
		IN p_inventario_id BIGINT,
		IN p_dias_prestamo INT DEFAULT 7,
		INOUT p_prestamo_id BIGINT DEFAULT NULL
		)
		LANGUAGE plpgsql
		AS $$
		DECLARE
		v_estado_inventario VARCHAR;
		v_prestamos_activos BIGINT;
		BEGIN
		SELECT estado INTO v_estado_inventario FROM inventario WHERE id = p_inventario_id;
		IF v_estado_inventario IS NULL THEN
		RAISE EXCEPTION 'Ejemplar no encontrado';
		END IF;
		IF v_estado_inventario != 'DISPONIBLE' THEN
		RAISE EXCEPTION 'El ejemplar no está disponible';
		END IF;
		
		SELECT COUNT(*) INTO v_prestamos_activos
		FROM prestamos WHERE usuario_id = p_usuario_id AND estado = 'ACTIVO';
		IF v_prestamos_activos >= 5 THEN
		RAISE EXCEPTION 'El usuario tiene demasiados préstamos activos';
		END IF;
		
		INSERT INTO prestamos (usuario_id, inventario_id, fecha_prestamo, fecha_vencimiento, estado)
		VALUES (p_usuario_id, p_inventario_id, NOW(), NOW() + (p_dias_prestamo || ' days')::INTERVAL, 'ACTIVO')
		RETURNING id INTO p_prestamo_id;
		
		UPDATE inventario SET estado = 'PRESTADO' WHERE id = p_inventario_id;
		
		UPDATE libros l
		SET ejemplares_disponibles = ejemplares_disponibles - 1
		FROM inventario i
		WHERE i.id = p_inventario_id AND l.id = i.libro_id;
		END;
		$$;
	\end{lstlisting}
	
	\hypertarget{sp_devolver_prestamo.sql}{%
		\section{sp\_devolver\_prestamo.sql}\label{sp_devolver_prestamo.sql}}
	
	\passthrough{\lstinline!db/procs/sp\_devolver\_prestamo.sql!}
	
	\begin{lstlisting}[language=SQL]
		CREATE OR REPLACE PROCEDURE sp_devolver_prestamo(
		IN p_prestamo_id BIGINT,
		INOUT p_multa_generada BOOLEAN DEFAULT FALSE
		)
		LANGUAGE plpgsql
		AS $$
		DECLARE
		v_estado VARCHAR;
		v_inventario_id BIGINT;
		v_libro_id BIGINT;
		v_fecha_vencimiento TIMESTAMP;
		v_dias_retraso INT;
		BEGIN
		SELECT estado, inventario_id, fecha_vencimiento
		INTO v_estado, v_inventario_id, v_fecha_vencimiento
		FROM prestamos WHERE id = p_prestamo_id;
		
		IF v_estado IS NULL THEN
		RAISE EXCEPTION 'Préstamo no encontrado';
		END IF;
		IF v_estado = 'DEVUELTO' THEN
		RAISE EXCEPTION 'El préstamo ya fue devuelto';
		END IF;
		
		UPDATE prestamos
		SET estado = 'DEVUELTO', fecha_devolucion = NOW()
		WHERE id = p_prestamo_id;
		
		UPDATE inventario SET estado = 'DISPONIBLE' WHERE id = v_inventario_id;
		
		UPDATE libros l
		SET ejemplares_disponibles = ejemplares_disponibles + 1
		FROM inventario i
		WHERE i.id = v_inventario_id AND l.id = i.libro_id;
		
		IF v_fecha_vencimiento < NOW() THEN
		v_dias_retraso := EXTRACT(DAY FROM (NOW() - v_fecha_vencimiento));
		INSERT INTO multas (prestamo_id, usuario_id, monto, pagada, fecha_multa)
		SELECT p_prestamo_id, p.usuario_id, v_dias_retraso * 0.50, false, NOW()
		FROM prestamos p WHERE p.id = p_prestamo_id;
		p_multa_generada := TRUE;
		END IF;
		END;
		$$;
	\end{lstlisting}
	
	\hypertarget{sp_reporte_prestamos_diarios.sql}{%
		\section{sp\_reporte\_prestamos\_diarios.sql}\label{sp_reporte_prestamos_diarios.sql}}
	
	\passthrough{\lstinline!db/procs/sp\_reporte\_prestamos\_diarios.sql!}
	
	\begin{lstlisting}[language=SQL]
		CREATE OR REPLACE PROCEDURE sp_reporte_prestamos_diarios(
		IN p_fecha DATE,
		INOUT p_ref REFCURSOR
		)
		LANGUAGE plpgsql
		AS $$
		BEGIN
		OPEN p_ref FOR
		SELECT p.id, p.fecha_prestamo, p.estado,
		u.nombre AS usuario_nombre,
		l.titulo AS libro_titulo
		FROM prestamos p
		JOIN usuarios u ON u.id = p.usuario_id
		JOIN inventario i ON i.id = p.inventario_id
		JOIN libros l ON l.id = i.libro_id
		WHERE p.fecha_prestamo::DATE = p_fecha
		ORDER BY p.fecha_prestamo;
		END;
		$$;
	\end{lstlisting}
	
	\hypertarget{apendices}{%
		\chapter{7. Apendices}\label{apendices}}
	
	\hypertarget{bituxe1cora-de-observaciones-sigcb-qr}{%
		\section{Bitácora de Observaciones --- SIGCB-QR}\label{bituxe1cora-de-observaciones-sigcb-qr}}
	
	\begin{quote}
		Archivo de trazabilidad de observaciones recibidas del docente en las Entregas 1A y 1B,
		conforme al requisito \textbf{Bloque 0} de la Tercera Entrega (v0.9.0-rc).
	\end{quote}
	
	\hypertarget{resumen-1}{%
		\subsection{Resumen}\label{resumen-1}}
	
	\begin{longtable}[]{@{}ll@{}}
		\toprule\noalign{}
		Indicador & Valor \\
		\midrule\noalign{}
		\endhead
		\bottomrule\noalign{}
		\endlastfoot
		Total observaciones & 7 \\
		Resueltas & 7 \\
		No resueltas & 0 \\
		\% Resuelto & 100 \% \\
	\end{longtable}
	
	\hypertarget{observaciones}{%
		\subsection{Observaciones}\label{observaciones}}
	
	\begin{longtable}[]{@{}
			>{\raggedright\arraybackslash}p{(\columnwidth - 10\tabcolsep) * \real{0.1311}}
			>{\raggedright\arraybackslash}p{(\columnwidth - 10\tabcolsep) * \real{0.1311}}
			>{\raggedright\arraybackslash}p{(\columnwidth - 10\tabcolsep) * \real{0.1639}}
			>{\raggedright\arraybackslash}p{(\columnwidth - 10\tabcolsep) * \real{0.2131}}
			>{\raggedright\arraybackslash}p{(\columnwidth - 10\tabcolsep) * \real{0.1639}}
			>{\raggedright\arraybackslash}p{(\columnwidth - 10\tabcolsep) * \real{0.1967}}@{}}
		\toprule\noalign{}
		\begin{minipage}[b]{\linewidth}\raggedright
			Código
		\end{minipage} & \begin{minipage}[b]{\linewidth}\raggedright
			Fuente
		\end{minipage} & \begin{minipage}[b]{\linewidth}\raggedright
			Criterio
		\end{minipage} & \begin{minipage}[b]{\linewidth}\raggedright
			Observación
		\end{minipage} & \begin{minipage}[b]{\linewidth}\raggedright
			Decisión
		\end{minipage} & \begin{minipage}[b]{\linewidth}\raggedright
			Commit(es)
		\end{minipage} \\
		\midrule\noalign{}
		\endhead
		\bottomrule\noalign{}
		\endlastfoot
		OBS-01 & Entrega 1A & SRS & El documento SRS no sigue la plantilla ISO/IEC/IEEE 29148:2018. Faltan secciones de rationale, trazabilidad y métodos de verificación. & Implementado: SRS completo en docs/requisitos/SRS.md con todas las secciones obligatorias, identificadores únicos, rationale, trazabilidad y criterios de verificación. & \\
		OBS-02 & Entrega 1B & Auth JWT & El JWT no incluye el claim \passthrough{\lstinline!aud!} (audience) para restringir el público objetivo del token. & Aplicada: se agregaron los claims \passthrough{\lstinline!iss!}, \passthrough{\lstinline!aud!}, \passthrough{\lstinline!nbf!} e \passthrough{\lstinline!iat!} al JWT firmado. & \\
		OBS-03 & Entrega 1B & Auth Cookie & La cookie JWT tiene \passthrough{\lstinline!Secure=false!} en desarrollo; debe obtenerse de configuración externa. & Aplicada: \passthrough{\lstinline!Secure!} ahora se configura desde \passthrough{\lstinline!app.jwt.secure-cookie!} en application.yml. & \\
		OBS-04 & Entrega 1B & API Errors & Los errores HTTP no siguen RFC 7807 (ProblemDetails). No se distingue entre type, title, status, detail e instance. & Aplicada: GlobalExceptionHandler ahora retorna \passthrough{\lstinline!ProblemDetail!} con los campos type, title, status, detail e instance. & \\
		OBS-05 & Entrega 1B & Cache & No hay cache Redis implementado. El endpoint de listado de libros no tiene TTL configurable externamente ni métrica de hit ratio. & Aplicada: Se integró Redis con TTL desde \passthrough{\lstinline!app.cache.default-ttl-seconds!}. Se agregó métrica de hit ratio en docs/mediciones/. & \\
		OBS-06 & Entrega 1B & Data Access & Las consultas de reportes y agregaciones se realizan vía JPQL en lugar de procedimientos almacenados. No hay separación CRUD-ORM vs SP. & Aplicada: Se migraron todas las consultas no elementales a stored procedures en \passthrough{\lstinline!db/procs/!}. Catálogo en docs/basedatos/CATALOGO-SP.md. & \\
		OBS-07 & Entrega 1B & Docker & Las imágenes Docker usan etiquetas variables (\passthrough{\lstinline!latest!}, \passthrough{\lstinline!16-alpine!}) en lugar de digest SHA256, lo que rompe la reproducibilidad. & Aplicada: Todas las imágenes fijadas por digest SHA256 en docker-compose.yml. & \\
	\end{longtable}
	
	\hypertarget{historial-de-cambios}{%
		\subsection{Historial de cambios}\label{historial-de-cambios}}
	
	\begin{longtable}[]{@{}
			>{\raggedright\arraybackslash}p{(\columnwidth - 4\tabcolsep) * \real{0.3182}}
			>{\raggedright\arraybackslash}p{(\columnwidth - 4\tabcolsep) * \real{0.3182}}
			>{\raggedright\arraybackslash}p{(\columnwidth - 4\tabcolsep) * \real{0.3636}}@{}}
		\toprule\noalign{}
		\begin{minipage}[b]{\linewidth}\raggedright
			Fecha
		\end{minipage} & \begin{minipage}[b]{\linewidth}\raggedright
			Autor
		\end{minipage} & \begin{minipage}[b]{\linewidth}\raggedright
			Cambio
		\end{minipage} \\
		\midrule\noalign{}
		\endhead
		\bottomrule\noalign{}
		\endlastfoot
		2026-07-30 & Equipo SIGCB-QR & Creación inicial del archivo con 7 observaciones identificadas del análisis técnico. Pendiente de completar con feedback formal del docente. \\
	\end{longtable}
	
	\hypertarget{changelog-de-requisitos-sigcb-qr}{%
		\section{Changelog de Requisitos --- SIGCB-QR}\label{changelog-de-requisitos-sigcb-qr}}
	
	\begin{quote}
		Formato: \href{https://keepachangelog.com/}{Keep a Changelog} adaptado a requisitos.
	\end{quote}
	
	\hypertarget{v0.9.0-rc-2026-07-30}{%
		\subsection{{[}v0.9.0-rc{]} --- 2026-07-30}\label{v0.9.0-rc-2026-07-30}}
	
	\hypertarget{added}{%
		\subsubsection{Added}\label{added}}
	
	\begin{itemize}
		\tightlist
		\item
		REQ-F-008: Renovacion de prestamo
		\item
		REQ-F-009: Reporte de prestamos diarios
		\item
		REQ-F-010: Dashboard de estadisticas
		\item
		REQ-F-011: Gestion de reservas
		\item
		REQ-NF-005: Cobertura de pruebas \textgreater= 30\% (JaCoCo)
		\item
		REQ-NF-006: Control de acceso por roles (RBAC)
		\item
		REQ-NF-007: Cache Redis para blacklist JWT
		\item
		REQ-R-001: Tiempo de respuesta de autenticacion (p95 \textless= 500ms)
		\item
		REQ-R-002: Tiempo de respuesta de CRUD (p95 \textless= 300ms)
		\item
		REQ-R-003: Tiempo de respuesta de stored procedures (\textless{} 1s)
		\item
		REQ-R-004: Hit ratio de cache Redis \textgreater= 80\%
		\item
		HU-06: Prestamo de ejemplares
		\item
		HU-07: Devolucion de prestamo
		\item
		HU-08: Renovacion de prestamo
		\item
		HU-09: Reportes bibliotecarios
		\item
		HU-10: Dashboard de estadisticas
		\item
		CU-03: Prestamo de ejemplares
		\item
		CU-04: Devolucion de prestamo
		\item
		CU-05: Generacion de reportes
		\item
		SRS estructurado completo segun ISO/IEC/IEEE 29148:2018
		\item
		Matriz de trazabilidad con 23 filas
		\item
		SRS.pdf generado
	\end{itemize}
	
	\hypertarget{modified}{%
		\subsubsection{Modified}\label{modified}}
	
	\begin{itemize}
		\tightlist
		\item
		REQ-F-001: Claims JWT ampliados (iss, aud, nbf, iat agregados)
		\item
		REQ-F-001: Cookie JWT ahora usa Secure desde configuracion externa
		\item
		REQ-F-003: Logout ahora usa blacklist JWT
		\item
		REQ-F-005: Cache Redis agregado con @Cacheable
		\item
		REQ-F-006: Migrado a stored procedure (sp\_crear\_prestamo)
		\item
		REQ-F-007: Migrado a stored procedure (sp\_devolver\_prestamo)
		\item
		REQ-NF-002: Endurecido con SameSite=Strict + Secure configurable
		\item
		Error handling: migrado de ApiResponse a RFC 7807 ProblemDetails
	\end{itemize}
	
	\hypertarget{v0.7.0-2026-07-14-entrega-1b}{%
		\subsection{{[}v0.7.0{]} --- 2026-07-14 (Entrega 1B)}\label{v0.7.0-2026-07-14-entrega-1b}}
	
	\hypertarget{added-1}{%
		\subsubsection{Added}\label{added-1}}
	
	\begin{itemize}
		\tightlist
		\item
		REQ-F-001: Autenticacion stateless JWT
		\item
		REQ-F-002: Registro de usuarios
		\item
		REQ-F-003: Cierre de sesion con JTI blacklist
		\item
		REQ-F-004: CRUD de usuarios (Admin)
		\item
		REQ-F-005: CRUD de material bibliografico
		\item
		REQ-F-006: Prestamo de ejemplares (via JPA)
		\item
		REQ-F-007: Devolucion de prestamo (via JPA)
		\item
		REQ-NF-001: Rendimiento de listado de libros
		\item
		REQ-NF-002: Seguridad de autenticacion
		\item
		REQ-NF-003: Proteccion contra inyeccion SQL
		\item
		REQ-NF-004: Disponibilidad mediante Docker Compose
	\end{itemize}
	
	\hypertarget{v0.3.0-2026-06-04-entrega-1a}{%
		\subsection{{[}v0.3.0{]} --- 2026-06-04 (Entrega 1A)}\label{v0.3.0-2026-06-04-entrega-1a}}
	
	\hypertarget{added-2}{%
		\subsubsection{Added}\label{added-2}}
	
	\begin{itemize}
		\tightlist
		\item
		SRS inicial con requisitos funcionales y no funcionales base
		\item
		Casos de uso y arquitectura C4 nivel 1 y 2
	\end{itemize}
	
	\hypertarget{catuxe1logo-de-procedimientos-almacenados-y-funciones-sigcb-qr}{%
		\section{Catálogo de Procedimientos Almacenados y Funciones --- SIGCB-QR}\label{catuxe1logo-de-procedimientos-almacenados-y-funciones-sigcb-qr}}
	
	\begin{quote}
		Versión: v0.9.0-rc
		Fecha: 2026-07-30
	\end{quote}
	
	\hypertarget{convenciuxf3n-de-nomenclatura}{%
		\subsection{Convención de nomenclatura}\label{convenciuxf3n-de-nomenclatura}}
	
	\begin{longtable}[]{@{}
			>{\raggedright\arraybackslash}p{(\columnwidth - 4\tabcolsep) * \real{0.3750}}
			>{\raggedright\arraybackslash}p{(\columnwidth - 4\tabcolsep) * \real{0.2500}}
			>{\raggedright\arraybackslash}p{(\columnwidth - 4\tabcolsep) * \real{0.3750}}@{}}
		\toprule\noalign{}
		\begin{minipage}[b]{\linewidth}\raggedright
			Prefijo
		\end{minipage} & \begin{minipage}[b]{\linewidth}\raggedright
			Tipo
		\end{minipage} & \begin{minipage}[b]{\linewidth}\raggedright
			Archivo
		\end{minipage} \\
		\midrule\noalign{}
		\endhead
		\bottomrule\noalign{}
		\endlastfoot
		\passthrough{\lstinline!fn\_!} & Función SQL (RETURNS) & \passthrough{\lstinline!db/procs/fn\_<verbo>\_<sustantivo>.sql!} \\
		\passthrough{\lstinline!sp\_!} & Procedimiento almacenado & \passthrough{\lstinline!db/procs/sp\_<verbo>\_<sustantivo>.sql!} \\
	\end{longtable}
	
	\hypertarget{funciones}{%
		\subsection{Funciones}\label{funciones}}
	
	\hypertarget{fn_contar_prestamos_entre_fechas}{%
		\subsubsection{fn\_contar\_prestamos\_entre\_fechas}\label{fn_contar_prestamos_entre_fechas}}
	
	\begin{longtable}[]{@{}ll@{}}
		\toprule\noalign{}
		Propiedad & Valor \\
		\midrule\noalign{}
		\endhead
		\bottomrule\noalign{}
		\endlastfoot
		Propósito & Cuenta préstamos en un rango de fechas \\
		Parámetros & \passthrough{\lstinline!p\_inicio TIMESTAMP!}, \passthrough{\lstinline!p\_fin TIMESTAMP!} \\
		Retorno & \passthrough{\lstinline!BIGINT!} \\
		Tablas afectadas & \passthrough{\lstinline!prestamos!} (SELECT) \\
	\end{longtable}
	
	\hypertarget{fn_total_multas_pendientes}{%
		\subsubsection{fn\_total\_multas\_pendientes}\label{fn_total_multas_pendientes}}
	
	\begin{longtable}[]{@{}ll@{}}
		\toprule\noalign{}
		Propiedad & Valor \\
		\midrule\noalign{}
		\endhead
		\bottomrule\noalign{}
		\endlastfoot
		Propósito & Suma el monto total de multas no pagadas \\
		Parámetros & Ninguno \\
		Retorno & \passthrough{\lstinline!DECIMAL(12,2)!} \\
		Tablas afectadas & \passthrough{\lstinline!multas!} (SELECT) \\
	\end{longtable}
	
	\hypertarget{fn_total_cobrado_entre}{%
		\subsubsection{fn\_total\_cobrado\_entre}\label{fn_total_cobrado_entre}}
	
	\begin{longtable}[]{@{}ll@{}}
		\toprule\noalign{}
		Propiedad & Valor \\
		\midrule\noalign{}
		\endhead
		\bottomrule\noalign{}
		\endlastfoot
		Propósito & Suma el monto cobrado en un rango de fechas \\
		Parámetros & \passthrough{\lstinline!p\_inicio TIMESTAMP!}, \passthrough{\lstinline!p\_fin TIMESTAMP!} \\
		Retorno & \passthrough{\lstinline!DECIMAL(12,2)!} \\
		Tablas afectadas & \passthrough{\lstinline!multas!} (SELECT) \\
	\end{longtable}
	
	\hypertarget{fn_libros_disponibles}{%
		\subsubsection{fn\_libros\_disponibles}\label{fn_libros_disponibles}}
	
	\begin{longtable}[]{@{}
			>{\raggedright\arraybackslash}p{(\columnwidth - 2\tabcolsep) * \real{0.6111}}
			>{\raggedright\arraybackslash}p{(\columnwidth - 2\tabcolsep) * \real{0.3889}}@{}}
		\toprule\noalign{}
		\begin{minipage}[b]{\linewidth}\raggedright
			Propiedad
		\end{minipage} & \begin{minipage}[b]{\linewidth}\raggedright
			Valor
		\end{minipage} \\
		\midrule\noalign{}
		\endhead
		\bottomrule\noalign{}
		\endlastfoot
		Propósito & Lista libros con ejemplares disponibles \\
		Parámetros & Ninguno \\
		Retorno & \passthrough{\lstinline!TABLE(id BIGINT, titulo VARCHAR, isbn VARCHAR, ejemplares\_disponibles INT)!} \\
		Tablas afectadas & \passthrough{\lstinline!libros!} (SELECT) \\
	\end{longtable}
	
	\hypertarget{procedimientos-almacenados}{%
		\subsection{Procedimientos almacenados}\label{procedimientos-almacenados}}
	
	\hypertarget{sp_reporte_libros_mas_prestados}{%
		\subsubsection{sp\_reporte\_libros\_mas\_prestados}\label{sp_reporte_libros_mas_prestados}}
	
	\begin{longtable}[]{@{}ll@{}}
		\toprule\noalign{}
		Propiedad & Valor \\
		\midrule\noalign{}
		\endhead
		\bottomrule\noalign{}
		\endlastfoot
		Propósito & Obtiene los libros más prestados (top N) \\
		Parámetros & \passthrough{\lstinline!p\_limit INT!} (IN), \passthrough{\lstinline!p\_ref REFCURSOR!} (INOUT) \\
		Cursores & Retorna cursor con id, titulo, isbn, veces\_prestado \\
		Tablas afectadas & \passthrough{\lstinline!libros!} (SELECT) \\
	\end{longtable}
	
	\hypertarget{sp_prestamos_entre_fechas}{%
		\subsubsection{sp\_prestamos\_entre\_fechas}\label{sp_prestamos_entre_fechas}}
	
	\begin{longtable}[]{@{}
			>{\raggedright\arraybackslash}p{(\columnwidth - 2\tabcolsep) * \real{0.6111}}
			>{\raggedright\arraybackslash}p{(\columnwidth - 2\tabcolsep) * \real{0.3889}}@{}}
		\toprule\noalign{}
		\begin{minipage}[b]{\linewidth}\raggedright
			Propiedad
		\end{minipage} & \begin{minipage}[b]{\linewidth}\raggedright
			Valor
		\end{minipage} \\
		\midrule\noalign{}
		\endhead
		\bottomrule\noalign{}
		\endlastfoot
		Propósito & Obtiene préstamos en un rango de fechas con datos de usuario y libro \\
		Parámetros & \passthrough{\lstinline!p\_inicio TIMESTAMP!} (IN), \passthrough{\lstinline!p\_fin TIMESTAMP!} (IN), \passthrough{\lstinline!p\_ref REFCURSOR!} (INOUT) \\
		Cursores & Retorna cursor con datos completos del préstamo + usuario + libro \\
		Tablas afectadas & \passthrough{\lstinline!prestamos!}, \passthrough{\lstinline!usuarios!}, \passthrough{\lstinline!inventario!}, \passthrough{\lstinline!libros!} (SELECT, JOIN) \\
	\end{longtable}
	
	\hypertarget{sp_top_usuarios_prestamos}{%
		\subsubsection{sp\_top\_usuarios\_prestamos}\label{sp_top_usuarios_prestamos}}
	
	\begin{longtable}[]{@{}ll@{}}
		\toprule\noalign{}
		Propiedad & Valor \\
		\midrule\noalign{}
		\endhead
		\bottomrule\noalign{}
		\endlastfoot
		Propósito & Obtiene los usuarios con más préstamos (top N) \\
		Parámetros & \passthrough{\lstinline!p\_limit INT!} (IN), \passthrough{\lstinline!p\_ref REFCURSOR!} (INOUT) \\
		Cursores & Retorna cursor con id, nombre, email, total\_prestamos \\
		Tablas afectadas & \passthrough{\lstinline!usuarios!}, \passthrough{\lstinline!prestamos!} (SELECT, JOIN, GROUP BY) \\
	\end{longtable}
	
	\hypertarget{sp_dashboard_estadisticas}{%
		\subsubsection{sp\_dashboard\_estadisticas}\label{sp_dashboard_estadisticas}}
	
	\begin{longtable}[]{@{}
			>{\raggedright\arraybackslash}p{(\columnwidth - 2\tabcolsep) * \real{0.6111}}
			>{\raggedright\arraybackslash}p{(\columnwidth - 2\tabcolsep) * \real{0.3889}}@{}}
		\toprule\noalign{}
		\begin{minipage}[b]{\linewidth}\raggedright
			Propiedad
		\end{minipage} & \begin{minipage}[b]{\linewidth}\raggedright
			Valor
		\end{minipage} \\
		\midrule\noalign{}
		\endhead
		\bottomrule\noalign{}
		\endlastfoot
		Propósito & Obtiene estadísticas agregadas para el dashboard \\
		Parámetros & \passthrough{\lstinline!p\_fecha\_inicio TIMESTAMP!} (IN), \passthrough{\lstinline!p\_fecha\_fin TIMESTAMP!} (IN), \passthrough{\lstinline!p\_ref REFCURSOR!} (INOUT) \\
		Cursores & Retorna cursor con 6 métricas agregadas \\
		Tablas afectadas & \passthrough{\lstinline!prestamos!}, \passthrough{\lstinline!libros!}, \passthrough{\lstinline!usuarios!}, \passthrough{\lstinline!reservas!}, \passthrough{\lstinline!multas!} (SELECT, COUNT, SUM) \\
	\end{longtable}
	
	\hypertarget{sp_crear_prestamo}{%
		\subsubsection{sp\_crear\_prestamo}\label{sp_crear_prestamo}}
	
	\begin{longtable}[]{@{}
			>{\raggedright\arraybackslash}p{(\columnwidth - 2\tabcolsep) * \real{0.6111}}
			>{\raggedright\arraybackslash}p{(\columnwidth - 2\tabcolsep) * \real{0.3889}}@{}}
		\toprule\noalign{}
		\begin{minipage}[b]{\linewidth}\raggedright
			Propiedad
		\end{minipage} & \begin{minipage}[b]{\linewidth}\raggedright
			Valor
		\end{minipage} \\
		\midrule\noalign{}
		\endhead
		\bottomrule\noalign{}
		\endlastfoot
		Propósito & Crea un préstamo con validaciones de disponibilidad y límite \\
		Parámetros & \passthrough{\lstinline!p\_usuario\_id BIGINT!} (IN), \passthrough{\lstinline!p\_inventario\_id BIGINT!} (IN), \passthrough{\lstinline!p\_dias\_prestamo INT DEFAULT 7!} (IN), \passthrough{\lstinline!p\_prestamo\_id BIGINT DEFAULT NULL!} (INOUT) \\
		Validaciones & Disponibilidad del ejemplar, límite de 5 préstamos activos \\
		Tablas afectadas & \passthrough{\lstinline!prestamos!} (INSERT), \passthrough{\lstinline!inventario!} (UPDATE), \passthrough{\lstinline!libros!} (UPDATE) \\
	\end{longtable}
	
	\hypertarget{sp_devolver_prestamo}{%
		\subsubsection{sp\_devolver\_prestamo}\label{sp_devolver_prestamo}}
	
	\begin{longtable}[]{@{}
			>{\raggedright\arraybackslash}p{(\columnwidth - 2\tabcolsep) * \real{0.6111}}
			>{\raggedright\arraybackslash}p{(\columnwidth - 2\tabcolsep) * \real{0.3889}}@{}}
		\toprule\noalign{}
		\begin{minipage}[b]{\linewidth}\raggedright
			Propiedad
		\end{minipage} & \begin{minipage}[b]{\linewidth}\raggedright
			Valor
		\end{minipage} \\
		\midrule\noalign{}
		\endhead
		\bottomrule\noalign{}
		\endlastfoot
		Propósito & Procesa la devolución de un préstamo y genera multa si aplica \\
		Parámetros & \passthrough{\lstinline!p\_prestamo\_id BIGINT!} (IN), \passthrough{\lstinline!p\_multa\_generada BOOLEAN DEFAULT FALSE!} (INOUT) \\
		Validaciones & Estado del préstamo, fecha de vencimiento \\
		Tablas afectadas & \passthrough{\lstinline!prestamos!} (UPDATE), \passthrough{\lstinline!inventario!} (UPDATE), \passthrough{\lstinline!libros!} (UPDATE), \passthrough{\lstinline!multas!} (INSERT) \\
	\end{longtable}
	
	\hypertarget{sp_renovar_prestamo}{%
		\subsubsection{sp\_renovar\_prestamo}\label{sp_renovar_prestamo}}
	
	\begin{longtable}[]{@{}
			>{\raggedright\arraybackslash}p{(\columnwidth - 2\tabcolsep) * \real{0.6111}}
			>{\raggedright\arraybackslash}p{(\columnwidth - 2\tabcolsep) * \real{0.3889}}@{}}
		\toprule\noalign{}
		\begin{minipage}[b]{\linewidth}\raggedright
			Propiedad
		\end{minipage} & \begin{minipage}[b]{\linewidth}\raggedright
			Valor
		\end{minipage} \\
		\midrule\noalign{}
		\endhead
		\bottomrule\noalign{}
		\endlastfoot
		Propósito & Renueva un préstamo activo creando uno nuevo \\
		Parámetros & \passthrough{\lstinline!p\_prestamo\_id BIGINT!} (IN), \passthrough{\lstinline!p\_dias\_renovacion INT DEFAULT 7!} (IN), \passthrough{\lstinline!p\_nuevo\_prestamo\_id BIGINT DEFAULT NULL!} (INOUT) \\
		Validaciones & Estado del préstamo original \\
		Tablas afectadas & \passthrough{\lstinline!prestamos!} (UPDATE, INSERT) \\
	\end{longtable}
	
	\hypertarget{sp_reporte_prestamos_diarios}{%
		\subsubsection{sp\_reporte\_prestamos\_diarios}\label{sp_reporte_prestamos_diarios}}
	
	\begin{longtable}[]{@{}
			>{\raggedright\arraybackslash}p{(\columnwidth - 2\tabcolsep) * \real{0.6111}}
			>{\raggedright\arraybackslash}p{(\columnwidth - 2\tabcolsep) * \real{0.3889}}@{}}
		\toprule\noalign{}
		\begin{minipage}[b]{\linewidth}\raggedright
			Propiedad
		\end{minipage} & \begin{minipage}[b]{\linewidth}\raggedright
			Valor
		\end{minipage} \\
		\midrule\noalign{}
		\endhead
		\bottomrule\noalign{}
		\endlastfoot
		Propósito & Genera reporte de préstamos de un día específico \\
		Parámetros & \passthrough{\lstinline!p\_fecha DATE!} (IN), \passthrough{\lstinline!p\_ref REFCURSOR!} (INOUT) \\
		Cursores & Retorna cursor con préstamos del día + usuario + libro \\
		Tablas afectadas & \passthrough{\lstinline!prestamos!}, \passthrough{\lstinline!usuarios!}, \passthrough{\lstinline!inventario!}, \passthrough{\lstinline!libros!} (SELECT, JOIN) \\
	\end{longtable}
	
	\hypertarget{sp_reporte_multas_cobradas}{%
		\subsubsection{sp\_reporte\_multas\_cobradas}\label{sp_reporte_multas_cobradas}}
	
	\begin{longtable}[]{@{}
			>{\raggedright\arraybackslash}p{(\columnwidth - 2\tabcolsep) * \real{0.6111}}
			>{\raggedright\arraybackslash}p{(\columnwidth - 2\tabcolsep) * \real{0.3889}}@{}}
		\toprule\noalign{}
		\begin{minipage}[b]{\linewidth}\raggedright
			Propiedad
		\end{minipage} & \begin{minipage}[b]{\linewidth}\raggedright
			Valor
		\end{minipage} \\
		\midrule\noalign{}
		\endhead
		\bottomrule\noalign{}
		\endlastfoot
		Propósito & Genera reporte de multas cobradas en un mes específico \\
		Parámetros & \passthrough{\lstinline!p\_mes INT!} (IN), \passthrough{\lstinline!p\_anio INT!} (IN), \passthrough{\lstinline!p\_ref REFCURSOR!} (INOUT) \\
		Cursores & Retorna cursor con mes, total\_cobrado, pendientes \\
		Tablas afectadas & \passthrough{\lstinline!multas!} (SELECT, COUNT, SUM) \\
	\end{longtable}
	
	\hypertarget{seguridad}{%
		\subsection{Seguridad}\label{seguridad}}
	
	Todos los procedimientos usan exclusivamente parámetros nombrados. No se construye SQL dinámico
	mediante concatenación (EXECUTE IMMEDIATE, sp\_executesql o equivalentes).
	
	La cuenta de base de datos de la aplicación tiene únicamente:
	- \passthrough{\lstinline!EXECUTE!} sobre estos procedimientos y funciones
	- \passthrough{\lstinline!SELECT!}, \passthrough{\lstinline!INSERT!}, \passthrough{\lstinline!UPDATE!}, \passthrough{\lstinline!DELETE!} sobre las tablas del dominio
	
\end{document}